% Options for packages loaded elsewhere
\PassOptionsToPackage{unicode}{hyperref}
\PassOptionsToPackage{hyphens}{url}
\PassOptionsToPackage{dvipsnames,svgnames,x11names}{xcolor}
%
\documentclass[
  11pt,
  a4paper,
]{article}
\usepackage{amsmath,amssymb}
\usepackage{iftex}
\ifPDFTeX
  \usepackage[T1]{fontenc}
  \usepackage[utf8]{inputenc}
  \usepackage{textcomp} % provide euro and other symbols
\else % if luatex or xetex
  \usepackage{unicode-math} % this also loads fontspec
  \defaultfontfeatures{Scale=MatchLowercase}
  \defaultfontfeatures[\rmfamily]{Ligatures=TeX,Scale=1}
\fi
\usepackage{lmodern}
\ifPDFTeX\else
  % xetex/luatex font selection
  \setmainfont[]{Noto Serif}
  \setsansfont[]{Inter}
  \setmonofont[]{DejaVu Sans Mono}
  \setmathfont[]{Asana Math}
\fi
% Use upquote if available, for straight quotes in verbatim environments
\IfFileExists{upquote.sty}{\usepackage{upquote}}{}
\IfFileExists{microtype.sty}{% use microtype if available
  \usepackage[]{microtype}
  \UseMicrotypeSet[protrusion]{basicmath} % disable protrusion for tt fonts
}{}
\makeatletter
\@ifundefined{KOMAClassName}{% if non-KOMA class
  \IfFileExists{parskip.sty}{%
    \usepackage{parskip}
  }{% else
    \setlength{\parindent}{0pt}
    \setlength{\parskip}{6pt plus 2pt minus 1pt}}
}{% if KOMA class
  \KOMAoptions{parskip=half}}
\makeatother
\usepackage{xcolor}
\usepackage[margin=24mm]{geometry}
\usepackage{color}
\usepackage{fancyvrb}
\newcommand{\VerbBar}{|}
\newcommand{\VERB}{\Verb[commandchars=\\\{\}]}
\DefineVerbatimEnvironment{Highlighting}{Verbatim}{commandchars=\\\{\}}
% Add ',fontsize=\small' for more characters per line
\newenvironment{Shaded}{}{}
\newcommand{\AlertTok}[1]{\textcolor[rgb]{1.00,0.00,0.00}{\textbf{#1}}}
\newcommand{\AnnotationTok}[1]{\textcolor[rgb]{0.38,0.63,0.69}{\textbf{\textit{#1}}}}
\newcommand{\AttributeTok}[1]{\textcolor[rgb]{0.49,0.56,0.16}{#1}}
\newcommand{\BaseNTok}[1]{\textcolor[rgb]{0.25,0.63,0.44}{#1}}
\newcommand{\BuiltInTok}[1]{\textcolor[rgb]{0.00,0.50,0.00}{#1}}
\newcommand{\CharTok}[1]{\textcolor[rgb]{0.25,0.44,0.63}{#1}}
\newcommand{\CommentTok}[1]{\textcolor[rgb]{0.38,0.63,0.69}{\textit{#1}}}
\newcommand{\CommentVarTok}[1]{\textcolor[rgb]{0.38,0.63,0.69}{\textbf{\textit{#1}}}}
\newcommand{\ConstantTok}[1]{\textcolor[rgb]{0.53,0.00,0.00}{#1}}
\newcommand{\ControlFlowTok}[1]{\textcolor[rgb]{0.00,0.44,0.13}{\textbf{#1}}}
\newcommand{\DataTypeTok}[1]{\textcolor[rgb]{0.56,0.13,0.00}{#1}}
\newcommand{\DecValTok}[1]{\textcolor[rgb]{0.25,0.63,0.44}{#1}}
\newcommand{\DocumentationTok}[1]{\textcolor[rgb]{0.73,0.13,0.13}{\textit{#1}}}
\newcommand{\ErrorTok}[1]{\textcolor[rgb]{1.00,0.00,0.00}{\textbf{#1}}}
\newcommand{\ExtensionTok}[1]{#1}
\newcommand{\FloatTok}[1]{\textcolor[rgb]{0.25,0.63,0.44}{#1}}
\newcommand{\FunctionTok}[1]{\textcolor[rgb]{0.02,0.16,0.49}{#1}}
\newcommand{\ImportTok}[1]{\textcolor[rgb]{0.00,0.50,0.00}{\textbf{#1}}}
\newcommand{\InformationTok}[1]{\textcolor[rgb]{0.38,0.63,0.69}{\textbf{\textit{#1}}}}
\newcommand{\KeywordTok}[1]{\textcolor[rgb]{0.00,0.44,0.13}{\textbf{#1}}}
\newcommand{\NormalTok}[1]{#1}
\newcommand{\OperatorTok}[1]{\textcolor[rgb]{0.40,0.40,0.40}{#1}}
\newcommand{\OtherTok}[1]{\textcolor[rgb]{0.00,0.44,0.13}{#1}}
\newcommand{\PreprocessorTok}[1]{\textcolor[rgb]{0.74,0.48,0.00}{#1}}
\newcommand{\RegionMarkerTok}[1]{#1}
\newcommand{\SpecialCharTok}[1]{\textcolor[rgb]{0.25,0.44,0.63}{#1}}
\newcommand{\SpecialStringTok}[1]{\textcolor[rgb]{0.73,0.40,0.53}{#1}}
\newcommand{\StringTok}[1]{\textcolor[rgb]{0.25,0.44,0.63}{#1}}
\newcommand{\VariableTok}[1]{\textcolor[rgb]{0.10,0.09,0.49}{#1}}
\newcommand{\VerbatimStringTok}[1]{\textcolor[rgb]{0.25,0.44,0.63}{#1}}
\newcommand{\WarningTok}[1]{\textcolor[rgb]{0.38,0.63,0.69}{\textbf{\textit{#1}}}}
\usepackage{longtable,booktabs,array}
\usepackage{calc} % for calculating minipage widths
% Correct order of tables after \paragraph or \subparagraph
\usepackage{etoolbox}
\makeatletter
\patchcmd\longtable{\par}{\if@noskipsec\mbox{}\fi\par}{}{}
\makeatother
% Allow footnotes in longtable head/foot
\IfFileExists{footnotehyper.sty}{\usepackage{footnotehyper}}{\usepackage{footnote}}
\makesavenoteenv{longtable}
\setlength{\emergencystretch}{3em} % prevent overfull lines
\providecommand{\tightlist}{%
  \setlength{\itemsep}{0pt}\setlength{\parskip}{0pt}}
\setcounter{secnumdepth}{-\maxdimen} % remove section numbering
\ifLuaTeX
\usepackage[bidi=basic]{babel}
\else
\usepackage[bidi=default]{babel}
\fi
\babelprovide[main,import]{american}
\ifPDFTeX
\else
\babelfont{rm}[]{Noto Serif}
\fi
% get rid of language-specific shorthands (see #6817):
\let\LanguageShortHands\languageshorthands
\def\languageshorthands#1{}
\usepackage{microtype}
\usepackage{xcolor}
\usepackage{fancyhdr}
\usepackage{titlesec}
\usepackage{etoolbox}
\usepackage{xurl}
\usepackage{needspace}
\definecolor{CSGNInk}{HTML}{203745}
\definecolor{CSGNTeal}{HTML}{1B6771}
\definecolor{CSGNMuted}{HTML}{667681}
\AtBeginDocument{\hypersetup{linkcolor=CSGNTeal,urlcolor=CSGNTeal,citecolor=CSGNTeal,pdftitle={CSGN v0.4.1: Objective-based plasticity, bounded graph computation, and verifiable consolidation},pdfauthor={NeuroPlastic Research Project}}}
\titleformat{\section}{\Large\sffamily\bfseries\color{CSGNInk}}{}{0pt}{}
\titleformat{\subsection}{\large\sffamily\bfseries\color{CSGNTeal}}{}{0pt}{}
\titlespacing*{\section}{0pt}{2.0ex plus .4ex}{0.8ex}
\titlespacing*{\subsection}{0pt}{1.5ex plus .3ex}{0.5ex}
\setlength{\parindent}{0pt}
\setlength{\parskip}{0.45em}
\linespread{1.05}
\setlength{\emergencystretch}{3em}
\setlength{\tabcolsep}{5pt}
\renewcommand{\arraystretch}{1.13}
\AtBeginEnvironment{longtable}{\small}
\AtBeginEnvironment{verbatim}{\small}
\pagestyle{fancy}
\fancyhf{}
\fancyhead[L]{\sffamily\scriptsize\color{CSGNMuted}NEUROPLASTIC / RESEARCH SPECIFICATION}
\fancyhead[R]{\sffamily\scriptsize\color{CSGNMuted}CSGN v0.4.1}
\fancyfoot[L]{\sffamily\scriptsize\color{CSGNMuted}16 September 2026}
\fancyfoot[R]{\sffamily\scriptsize\color{CSGNMuted}\thepage}
\renewcommand{\headrulewidth}{0.3pt}
\setlength{\headheight}{14pt}
\setlength{\footskip}{24pt}
\preto\section{\Needspace{7\baselineskip}}
\preto\subsection{\Needspace{4\baselineskip}}
\widowpenalty=10000
\clubpenalty=10000
\displaywidowpenalty=10000
\allowdisplaybreaks[1]
\renewcommand{\contentsname}{Contents}
\renewcommand{\maketitle}{%
\begin{titlepage}
\sffamily\raggedright
{\color{CSGNTeal}\small NEUROPLASTIC RESEARCH PROJECT}\par
\vspace{28mm}
{\fontsize{32}{38}\selectfont\bfseries\color{CSGNInk}Cortical Synaptic\par Graph Networks}\par
\vspace{8mm}
{\fontsize{17}{23}\selectfont\color{CSGNInk}Objective-based plasticity,\par bounded graph computation,\par and verifiable consolidation}\par
\vspace{13mm}
{\large\bfseries\color{CSGNTeal}VERSION 0.4.1}\par
\vspace{2mm}
{\normalsize 16 September 2026}\par
\vspace{16mm}
{\color{CSGNTeal}\rule{\textwidth}{1.1pt}}\par
\vspace{3mm}
{\normalsize Prediction\ \(\rightarrow\)\ Error\ \(\rightarrow\)\ Bounded write\ \(\rightarrow\)\ Cold recall}\par
\vfill
{\small\bfseries Research specification, not an empirical validation.}\par
\vspace{2mm}
{\small Local guarantees are stated with their assumptions.\par Integrated learning, retention, robustness, and scale remain to be tested.}\par
\vspace{7mm}
{\color{CSGNMuted}\small Includes derivations, a lifecycle protocol, an experiment plan,\par and reproducible local mathematical checks.}\par
\end{titlepage}
}
\ifLuaTeX
  \usepackage{selnolig}  % disable illegal ligatures
\fi
\usepackage{bookmark}
\IfFileExists{xurl.sty}{\usepackage{xurl}}{} % add URL line breaks if available
\urlstyle{same}
\hypersetup{
  pdftitle={Cortical Synaptic Graph Networks},
  pdfauthor={NeuroPlastic Research Project},
  pdflang={en-US},
  colorlinks=true,
  linkcolor={Maroon},
  filecolor={Maroon},
  citecolor={Blue},
  urlcolor={Blue},
  pdfcreator={LaTeX via pandoc}}

\title{Cortical Synaptic Graph Networks}
\usepackage{etoolbox}
\makeatletter
\providecommand{\subtitle}[1]{% add subtitle to \maketitle
  \apptocmd{\@title}{\par {\large #1 \par}}{}{}
}
\makeatother
\subtitle{Objective-based plasticity, bounded graph computation, and
verifiable consolidation}
\author{NeuroPlastic Research Project}
\date{Version 0.4.1 · 16 September 2026}

\begin{document}
\maketitle

{
\hypersetup{linkcolor=}
\setcounter{tocdepth}{2}
\tableofcontents
}
\clearpage
\section*{Research status and
abstract}\label{research-status-and-abstract}
\addcontentsline{toc}{section}{Research status and abstract}

\textbf{Status: research specification with local mathematical
guarantees. The integrated architecture has not been empirically
validated.} This is a new manuscript, not a report of a trained agent or
a claim that the complete system converges. It supersedes v0.4 as the
authoritative proposed specification. Earlier releases remain historical
evidence, not alternate implementation defaults. The existing
implementation is outside its scope.

Cortical Synaptic Graph Networks (CSGN) are proposed sparse recurrent
learners with a fixed synapse-slot budget, rapidly writable residual
weights, slower consolidated weights, finite working and episodic
memory, and optional structural adaptation. The v0.4 line changes the
organizing principle from bounded correlation accumulation to
\textbf{bounded correction of an explicit local error}. Computation,
environmental interaction, plastic writes, and consolidation occur on
distinct clocks. During each write, its feature snapshot, routing
support, target, and feasible set are fixed. This permits exact
current-association error identities and projected-gradient descent
guarantees for a declared local objective. It does not imply improvement
of an arbitrary downstream policy.

The graph uses context-dependent routing with a null route, typed
bounded messages, and receiver budgets independent of the writable
weights. A sufficient contraction condition is given for the
computational workspace during a frozen internal cycle; persistent
storage is not forced to contract. Retention and writing are
independently controlled. Slow consolidation is evaluated with plastic
residuals cleared and transient memory standardized. Instantaneous
fast-to-slow transfer, long-term persistence, quantization error, and
behavioral improvement are treated as different claims.

The specification includes provenance-aware memory handling, a
synchronous reference commit protocol, bounded structural proposals,
explicit resource accounting, and a staged evaluation plan. Twenty-four
groups of newly executed algebra, feasibility, and gradient checks
accompany this revision. They validate local identities and demonstrate
counterexamples, not the capabilities of a complete CSGN. The main
empirical hypothesis is that learned representations and targeted
plastic writes can improve the adaptation--retention trade-off beyond
matched fixed-graph, replay-only, and established fast-memory baselines.
Success on ideal orthogonal keys is not sufficient evidence for that
hypothesis.

\textbf{Revision 0.4.1.} This revision replaces the overly conservative
absolute-sum reserve with exact four-endpoint decay-safe feasibility;
separates fixed inner variables from retained outer gradient paths;
makes target reachability and downstream write utility diagnostic
requirements; adds memory transplantation, selective-erasure,
negative-control, and rule-transfer experiments; elevates the simple
delta-memory comparator; and separates practical empirical validation
from formal statistical claims. These changes refine a small complete
lifecycle rather than add a large subsystem.

\textbf{Central design contract.} A memory write must identify what it
predicts, what evidence supplies its target, which state it changes, and
which loss its update can improve. Sleep must demonstrate retained
behavior after fast memory is removed. Scale follows those results
rather than substituting for them.

\section{1. Scope, motivation, and relationship to prior
work}\label{scope-motivation-and-relationship-to-prior-work}

\subsection{1.1 The research question}\label{the-research-question}

Can an agent learn new associations during operation, preserve useful
older associations, revise obsolete ones, and consolidate selected
knowledge into slower storage while using bounded memory and
computation? CSGN tests whether a sparse synaptic graph is a useful
substrate for that combination.

The proposal does not require a pretrained language model. It does
require a training process that makes perception, context
representations, routing, and local teaching signals useful. A randomly
initialized self-modifying circuit is not presumed to acquire these
capabilities merely by operating for a long time. Biological terminology
describes architectural inspiration, not an equivalence between
artificial columns, parameters, and biological neurons.

Three hypotheses are distinguished:

\textbf{H1 --- useful writing.} Objective-based plastic updates improve
held-out adaptation when representations and routing are learned, not
supplied as ideal keys.

\textbf{H2 --- useful consolidation.} Sleep improves slow-memory recall
after residuals and unrelated transient cues are removed, without
unacceptable losses on still-valid prior tasks.

\textbf{H3 --- useful structure.} Budgeted rewiring improves performance
or efficiency over a trained fixed graph after charging for its search,
retraining, validation, and data movement.

H3 is optional for the first demonstration. A complete small lifecycle
must still include graph computation, writable synapses, slow memory,
consolidation, fast-state clearing, and delayed recall. Component
ablations isolate causes; they do not replace the integrated experiment.

\subsection{1.2 Established ingredients and proposed
contribution}\label{established-ingredients-and-proposed-contribution}

Differentiable plasticity demonstrates that gradient-based training can
make Hebbian plastic connections useful on specific task families
{[}1{]}. Delta-rule fast-weight memories already provide targeted
corrections of key-to-value mappings {[}2{]}, and gated delta networks
combine retention control with such corrections {[}3{]}.
Test-time-training layers interpret recurrent state updates as
optimization of a learning objective {[}4{]}. CSGN does not claim
invention of these mechanisms.

Derivative-based eligibility and policy-gradient methods provide
precedents for associating parameter-specific sensitivities with
learning signals {[}5,6{]}. Dynamic sparse training supplies a reference
for fixed-budget topology adaptation {[}7{]}. Importance-based methods
motivate retention controls, but they are not guarantees of unlimited
lifelong memory {[}8{]}. Repeated adaptive validation requires its own
statistical discipline {[}9{]}, and low-precision training results do
not automatically establish safe INT4 lifetime learning {[}10{]}.

The proposed contribution is the \textbf{integration and evaluation} of
these ideas in a typed sparse graph, with an explicit write objective,
phase-specific guarantees, residual-aware consolidation, and honest
resource accounting. The local mathematical results below are derived
for this specification. Published results cited above are not evidence
that the complete CSGN works.

\subsection{1.3 Four meanings of
``works''}\label{four-meanings-of-works}

A well-defined update, a loss-reducing local update, a useful trained
learner, and an efficient lifetime agent are different milestones.
Version 0.4.1 specifies the first two under stated assumptions. It
proposes experiments for the third and fourth. None of the local
propositions establishes general intelligence, complete poisoning
resistance, universal non-forgetting, or practical operation at 100
billion stored edges.

\section{2. State, notation, and causal
clocks}\label{state-notation-and-causal-clocks}

\subsection{2.1 Persistent and transient
state}\label{persistent-and-transient-state}

There are \(N\) columns, each with \(k_{\mathrm{out}}\) reserved
outgoing slots. A slot \(e=(i,s)\) has source \(i\), target \(j(e)\),
message type \(\chi(e)\), structural class, and semantic generation
identifier. Targets may change only through a validated structural
transaction. Dormant slots remain part of the budget:

\[
E=N k_{\mathrm{out}},\qquad
k_{\mathrm{out}}=k_{\mathrm{backbone}}+k_{\mathrm{plastic}}.
\tag{1}
\]

The adaptive state at an external event boundary is

\[
\mathcal Z_t=
(\theta,S,P^a,P^u,h,\mathcal W,\mathcal M,
 G,\mathcal F,\mathcal V)_t.
\tag{2}
\]

Here \(\theta\) contains shared encoders, message transforms, routing,
workspace, and readout parameters; \(S\) is slow synaptic storage;
\(P^a\) and \(P^u\) are directly admissible and ephemeral residuals;
\(h\) is the recurrent workspace; \(\mathcal W\) is working memory;
\(\mathcal M\) is episodic memory; \(G\) is the versioned slot topology;
\(\mathcal F\) contains pending feedback records; and \(\mathcal V\)
contains model versions, quantization scales, timestamps, and controller
state.

The effective scalar synapse is

\[
w_e=S_e+P_e^a+P_e^u.
\tag{3}
\]

Residuals are in \textbf{effective-weight units}. The v0.3 expression
coefficient \(\alpha\) is not a separate online state; its scale is
absorbed into the residual parameterization. Write rates are controlled
separately. This is a redesign of the dynamics, not a claim that all
v0.3 trajectories are algebraically equivalent.

Cold slots have no active residual unless a retained residual record is
explicitly retrieved. The hot set \(\mathcal H_t\) includes resident
nonzero residuals and pinned pending-credit records. Discarding a
nonzero residual is a deliberate memory change, not a transparent cache
operation.

\subsection{2.2 Exact endpoint feasibility and independent
decay}\label{exact-endpoint-feasibility-and-independent-decay}

For every slot, require

\[
\begin{aligned}
\max_{a,b\in\{0,1\}}|S_e+aP_e^a+bP_e^u|&\le w_{\max},\\
|P_e^a|+|P_e^u|&\le p_{\max},
\end{aligned}
\qquad w_{\max}>0,\quad p_{\max}>0.
\tag{4}
\]

The four tested endpoints are \(S_e\), \(S_e+P_e^a\), \(S_e+P_e^u\), and
\(S_e+P_e^a+P_e^u\). This replaces v0.4's stronger condition
\(|S_e|+|P_e^a|+|P_e^u|\le w_{\max}\). The residual budget is retained
to bound internal cancellation and hot-state use.

\textbf{Endpoint lemma --- exact independent-decay safety.} For
\(\lambda_a,\lambda_u\in[0,1]\), the decayed value is

\[
\begin{aligned}
S+\lambda_aP^a+\lambda_uP^u={}&(1-\lambda_a)(1-\lambda_u)S\\
&+\lambda_a(1-\lambda_u)(S+P^a)\\
&+(1-\lambda_a)\lambda_u(S+P^u)\\
&+\lambda_a\lambda_u(S+P^a+P^u).
\end{aligned}
\]

The coefficients are nonnegative and sum to one, so the four endpoint
bounds are sufficient. They are also necessary for safety under
\emph{all} independent decays because each endpoint is attained by a
decay pair in \(\{0,1\}^2\). Applying additional independent decay to
the new residuals simply selects points in the original decay rectangle;
feasibility is therefore preserved under repeated decay. The residual
absolute-sum budget can only decrease.

\textbf{Corrective capacity.} At \(S=1\), \(w_{\max}=1\), \(P^u=0\), the
earlier reserve permits only \(P^a=0\). Equation (4) permits an inward
correction such as \(P^a=-0.3\) when \(p_{\max}\ge0.3\); all
independently decayed values remain between \(0.7\) and \(1\). This is
extra feasible corrective capacity, not a guarantee of better learned
behavior.

The feasible set is convex: the endpoint bounds are linear inequalities,
and the residual budget is an absolute-sum inequality. Workspace states
satisfy \(h_i\in[-1,1]^d\). Keys, values, contexts, and all arithmetic
inputs must be finite. Quantizer ranges, target ranges, and tolerances
are declared per experiment. Invalid starting states are rejected or
repaired as a separately logged maintenance operation, never silently
treated as a valid projection base.

\subsection{2.3 Four clocks and one event
order}\label{four-clocks-and-one-event-order}

Use \(t\) for an environmental decision, \(\ell\) for an internal
computation round, \(n\) for a plastic-write event, and \(\nu\) for a
maintenance transaction. The reference causal order is:

\begin{enumerate}
\def\labelenumi{\arabic{enumi}.}
\tightlist
\item
  Receive \(o_t\) and any feedback already available for earlier
  decisions. Resolve due feedback before the next decision, using the
  corresponding stored records.
\item
  Build bounded context from \(o_t\), prior working memory, and
  permitted retrieved records. Choose the active subgraph and freeze its
  route support, gates, slow weights, and residuals for this internal
  cycle.
\item
  Perform \(K_{\mathrm{inner}}\) workspace rounds. Produce the
  prediction or action. Save the write snapshot needed if a target will
  arrive later.
\item
  Receive the resulting observation, reward, or label when the
  environment makes it available. Create a causally valid target record;
  it cannot modify a decision that preceded that feedback.
\item
  Apply allowed decay and bounded writes. Maintenance occurs at a
  declared safe boundary.
\end{enumerate}

The reference online phase freezes \(\theta\). Offline/meta-training
updates \(\theta\) across episodes. A deployment that changes encoders
or shared transforms during maintenance must rebuild affected contexts
and records under the new version; it is an extension, not an invisible
implementation detail.

\section{3. Context-dependent sparse graph
routing}\label{context-dependent-sparse-graph-routing}

\subsection{3.1 Region and column
selection}\label{region-and-column-selection}

Partition columns into regions. A bounded context can be constructed as

\[
c_t=C_{\max}\tanh\!\left(
C_o x_t+C_w\,\operatorname{Read}(\mathcal W_{t^-})
+C_m\,\operatorname{Retrieve}(\mathcal M_t)
\right).
\tag{5}
\]

Retrieval is capped by an explicit record and byte budget. Disable it
during pure synaptic-recall tests. Its key encoder is frozen within an
online epoch and versioned across maintenance.

Select regions and then columns using query--key interactions. A local
score is

\[
a_i^t=
\frac{q_\theta(c_t)^\top k_\theta(h_i^{t,0},d_i)}{\sqrt{d_r}}
+b_i^{\mathrm{type}},
\tag{6}
\]

where \(d_i\) is a fixed or offline-learned column descriptor.
Descriptors break otherwise identical initial states. Query and key
outputs are normalized or norm-bounded. A shared additive context term
would cancel from the ordering and is not used.

For the small reference model, region scores and within-selected-region
scores are evaluated exactly. This costs \(O(Rd_r)\) plus the scored
local columns; it is not mislabeled as sublinear. Region summaries are
maintained incrementally when their member states change. At large
scale, replacing exact region scoring with a tree or approximate index
is a separately measured extension, with search recall, staleness, and
index-maintenance costs reported.

Denote active senders by \(\mathcal A_t\) and updated columns by
\(\mathcal U_t\), including admitted receivers and interface columns.
Absolute region, sender, receiver, and edge budgets are primary. The
active fraction is a resulting measurement, not an independent free
parameter.

\subsection{3.2 Gates, null mass, and
availability}\label{gates-null-mass-and-availability}

For a slot \(e=(i,s)\), calculate its routing logit from the current
context, source state, target descriptor/state snapshot, and message
type. Region-pair biases must be factorized or stored only for actual
regional relations; a dense \(R\times R\) table is not a scalable
default.

Let \(v_e\in[0,1]\) be structural availability. Exclude \(v_e=0\) from
selection. Form a bounded eligible support \(\mathcal S_i\) using
adjusted logits \(\ell_e+\log v_e\). Preserve a null route with finite
logit \(\ell_{i,0}\):

\[
\begin{aligned}
g_e&=\frac{v_e\exp\ell_e}{\exp\ell_{i,0}
+\sum_{u\in\mathcal S_i}v_u\exp\ell_u},\quad e\in\mathcal S_i,\\
g_{i,0}&=\frac{\exp\ell_{i,0}}{\exp\ell_{i,0}
+\sum_{u\in\mathcal S_i}v_u\exp\ell_u}.
\end{aligned}
\tag{7}
\]

Other slot gates are zero. Evaluate the expression with a stable
log-sum-exp calculation. Real routes satisfy \(\sum_s g_{i,s}\le1\). Do
not normalize them to total mass one and do not multiply by availability
again.

A deterministic receiver-capacity admission mask \(b_e\in\{0,1\}\) may
reject routes before computation. Define \(a_e=b_e g_e\ge0\) and freeze
it for the internal cycle and its later write snapshot. Rejected mass is
unused rather than silently redistributed. The set of scored candidate
edges, executed edges, and subsequently writable edges must be recorded
separately.

\textbf{Zero-availability property.} With unchanged selection of
incumbents, a new slot with \(v_e=0\) changes neither their gates nor
its own transmitted value. This does not make a later hard top-\(k\)
support switch continuous. Positive-availability admission and pruning
are evaluated as whole structural edits.

\subsection{3.3 Temporal service, not merely static
connectivity}\label{temporal-service-not-merely-static-connectivity}

A protected backbone supplies structural routes between regions, but a
static path is not necessarily executed in useful temporal order. Demand
activation queues should promote relevant receivers into future sender
sets and record time-respecting path latency. The union of edges across
a window is insufficient evidence of temporal reachability.

The reference scheduler has a finite demand queue, maximum admitted
work, queue overflow policy, and a maximum internal computation budget.
An overflow can defer a request, retain it externally, or reject a
plastic write; it must be observable in metrics. If at most \(K_R\)
distinct regions are activated per step, visiting all \(R\) regions
requires at least \(\lceil R/K_R\rceil\) steps. No constant-latency
full-coverage claim is made as storage grows.

\section{4. Bounded messages and a conditionally stable
workspace}\label{bounded-messages-and-a-conditionally-stable-workspace}

\subsection{4.1 Typed messages and weight-independent receiver
budgets}\label{typed-messages-and-weight-independent-receiver-budgets}

For each active sender and required head, use

\[
m_{i,q}^{t,\ell}
=M_{\max}\tanh(U_q h_i^{t,\ell}+V_q c_t),
\qquad q\in\{1,\ldots,H\}.
\tag{8}
\]

Thus \(\|m_{i,q}\|_2\le M=M_{\max}\sqrt{d_m}\). With fixed context, its
hidden-state Lipschitz constant is bounded by
\(L_m=M_{\max}\max_q\|U_q\|_2\). More expressive bounded transforms are
extensions requiring their own constant or measured stability regime.

Receiver \(j\) uses an envelope independent of the writable values:

\[
Z_j=\max\!\left(1,w_{\max}\sum_{e\to j}a_e\right),
\qquad
r_j^{t,\ell}=\frac1{Z_j}\sum_{e\to j}a_e w_e m_{i(e),\chi(e)}^{t,\ell}.
\tag{9}
\]

\textbf{Proposition 1 --- receiver bound.} Under Equation (4) and the
message bound,

\[
\|r_j\|_2
\le\frac{M\sum_{e\to j}a_e|w_e|}{Z_j}
\le M.
\tag{10}
\]

The first inequality is the triangle inequality; the second uses
\(|w_e|\le w_{\max}\) and the definition of \(Z_j\). The bound does not
depend on in-degree. This envelope can attenuate weak weights more than
a realized-weight normalizer; its signal scale is an experimental
trade-off, not a claim of universal superiority.

\subsection{4.2 Inner recurrent
computation}\label{inner-recurrent-computation}

For \(j\in\mathcal U_t\), use a fixed mixing coefficient
\(0<\kappa\le1\) during the cycle:

\[
h_j^{t,\ell+1}
=(1-\kappa)h_j^{t,\ell}
+\kappa\tanh(A_hh_j^{t,\ell}+B_rr_j^{t,\ell}+u_{j,t}),
\tag{11}
\]

where \(u_{j,t}\) is a bounded sensory/context/working-memory drive
fixed during the cycle. Unupdated columns retain their state. Any
declared decay of those states is applied at an event boundary and
preserves their box bound. A readout query is derived from \(c_t\),
rather than being an unspecified constant; the head predicts from
selected final states and permitted memory.

\textbf{Proposition 2 --- box invariance.} With initial
\(h_j\in[-1,1]^d\), Equation (11) preserves that set. Both endpoints of
its coordinatewise convex combination lie in the box. This is a
magnitude bound, not a claim about learned behavior.

\textbf{Proposition 3 --- frozen-cycle contraction.} Let \(a=\|A_h\|_2\)
and \(b=\|B_r\|_2\). In the block-maximum norm
\(\|h\|_{\mathrm{bm}}=\max_j\|h_j\|_2\), if

\[
a+bL_m<1,
\qquad q=1-\kappa+\kappa(a+bL_m)<1,
\tag{12}
\]

the update restricted to \(\mathcal U_t\) is contractive for fixed
gates, weights, input, context, and identical inactive-column boundary
states.

\emph{Proof.} The coefficient sum at a receiver satisfies
\(\sum_e a_e|w_e|/Z_j\le1\). Consequently its message change is at most
\(L_m\|\Delta h\|_{\mathrm{bm}}\). Tanh is 1-Lipschitz, so the candidate
update changes by at most \((a+bL_m)\|\Delta h\|_{\mathrm{bm}}\).
Combining the residual path gives the factor \(q\).

The reference profile enforces this condition by constraining operator
norms. For example, \(a\le0.25\), \(bL_m\le0.5\), and \(\kappa=0.5\)
give \(q\le0.875\); these are illustrative sufficient values, not tuned
recommendations. A less restrictive workspace is allowed only as a
separately labeled empirical variant. Sweep gains \emph{within} the
sufficient condition before relaxing it; excessive contraction can
attenuate useful signals. Record actual operator norms and task
performance. An approximate power-iteration estimate is not a certified
norm bound unless its estimation error is addressed.

\textbf{Scope limit.} The proposition is not contraction of all
persistent memory, the environment, routing switches, or the joint
plasticity loop. Frozen inactive coordinates have identity dynamics.
Between cycles, inputs and memory can change. Test finite-horizon
sensitivity of the augmented state as well as inner-cycle norms; stable
diagonal blocks do not imply stable coupled dynamics.

\section{5. Define what a plastic write
learns}\label{define-what-a-plastic-write-learns}

\subsection{5.1 Required target
contract}\label{required-target-contract}

Every write record contains a decision identifier, time of prediction,
time feedback became available, model/topology/feature versions, feature
snapshot, target, provenance class, and finite resource charge. The
target must be specified before invoking any learning guarantee.

The \textbf{reference association profile} predicts an externally
supplied value code from an observable cue and context. The value
arrives after the prediction. The predicted vector is a designated
receiver message, or a fixed linear decoding of selected messages. Use a
squared-error objective on that vector; classification by nearest code
is a downstream metric, not itself covered by the squared-error theorem.

The \textbf{predictive profile} uses a bounded encoding of a
subsequently observed quantity. Its target encoder is frozen within the
epoch and does not include the prediction being evaluated. Frozen
encoders stabilize coordinates but do not guarantee informative targets
or prevent all representation collapse.

A \textbf{learned local teacher} for internal receivers is an extension.
Its target is trained using task-grounded outer losses and compared
against full-gradient references. Decreasing its local loss is not
assumed to decrease final action loss. A novelty score, a confidence
score, or the network's own unsupported prediction is not an
independently verified target.

\subsection{5.2 A fixed-feature objective for scalar graph
synapses}\label{a-fixed-feature-objective-for-scalar-graph-synapses}

Capture the actual source messages, gates, receiver budgets, and
selected writable slot identities for a decision. They are constants
during the write. Group writable incoming edges into a vector \(z\) and
construct columns

\[
[A_n]_{j,e}=\frac{a_e}{Z_j}m_{i(e),\chi(e)}^{\mathrm{snapshot}}
\quad\text{in receiver }j\text{'s output coordinates}.
\tag{13}
\]

Columns are zero outside their receiver block. Subtract the contribution
of nonwritable edges, using their post-decay values held fixed for this
update, from the target to obtain \(b_n\). A fixed linear output decoder
can be absorbed into \(A_n\) and \(b_n\). Then

\[
\widehat b_n=A_nz,
\qquad
\ell_n(z)=\frac12\|A_nz-b_n\|_2^2,
\qquad
\nabla\ell_n=A_n^\top(A_nz-b_n).
\tag{14}
\]

This is a convex quadratic in the selected \textbf{graph synapses
themselves}. It does not add a mandatory external lookup model. The
weight-independent normalization in Equation (9) is essential to this
linear snapshot representation.

The snapshot is \textbf{numerically fixed for the partial derivative
with respect to writable synapses}. It is not automatically detached
from the outer meta-gradient. During exact short-episode meta-training,
retain the continuous dependency of the captured messages, gates, and
features on shared parameters; see Section 8.3. During ordinary
deployment, those shared parameters are frozen and snapshots need no
outer tape. On the next graph computation, hidden states and messages
may change, so the theorem below applies to the captured objective, not
automatically to the new unrolled network. With delayed feedback it is
an objective on the old snapshot evaluated at the current writable
values. Version-incompatible records must be reconstructed from
admissible experience or rejected, not applied to a reused slot address.

\subsection{5.3 Independent retention and
writing}\label{independent-retention-and-writing}

At write event \(n\), apply declared decay to residual tiers:

\[
P^{a,-}=\lambda_n^aP^a,\qquad
P^{u,-}=\lambda_n^uP^u,\qquad
0\le\lambda_n^a,\lambda_n^u\le1.
\tag{15}
\]

Persistent admissible associations use \(\lambda_n^a=1\) in the
reference retention test. Ephemeral decay may use
\(\exp(-\Delta t/\tau_u)\) with \(\tau_u>0\). Decay and learning have
separate budgets and logs. Equation (4) remains feasible after
independent decay.

The write strength \(\beta_n\in[0,1]\) is independent of \(\lambda_n\).
A memory may therefore receive a large event-driven write and retain it
without slow EMA attenuation. A lost association due to intentional
decay is not counted as a failure of the subsequent local-descent
proposition; all such comparisons explicitly start from the post-decay
state.

\section{6. Error-correcting updates and their
guarantees}\label{error-correcting-updates-and-their-guarantees}

\subsection{6.1 Normalized delta rule as a transparent special
case}\label{normalized-delta-rule-as-a-transparent-special-case}

Consider an allocated matrix block \(W=S+P\), a fixed nonzero key
\(k\in\mathbb R^{d_k}\), and a fixed target \(v\in\mathbb R^{d_v}\).
This block is a useful special case of synaptic memory and a diagnostic
benchmark; it need not replace the graph. After optional decay, let
\(W^-=S+\lambda P\) and \(e=v-W^-k\). Propose

\[
W^*=W^-+\frac{\beta}{\varepsilon+\|k\|_2^2}ek^\top,
\qquad 0\le\beta\le1,\quad\varepsilon\ge0.
\tag{16}
\]

Skip a zero key; it provides no write direction. The reference
finite-precision implementation uses a positive denominator floor. The
exact one-shot case below uses a nonzero key and \(\varepsilon=0\) only
to state the ideal identity.

\textbf{Proposition 4 --- current-association correction.} With fixed
\(k,v\) and no active projection or quantization,

\[
v-W^*k=
\left(1-\frac{\beta\|k\|_2^2}{\varepsilon+\|k\|_2^2}\right)e.
\tag{17}
\]

\emph{Proof.} Substitute Equation (16), multiply the outer product by
\(k\), and use \(k^\top k=\|k\|_2^2\). The scalar multiplier belongs to
\([0,1]\), so \(\tfrac12\|v-W^*k\|^2\le\tfrac12\|e\|^2\). For
\(\beta=1\) and \(\varepsilon=0\), the residual is zero in real
arithmetic.

This is the established delta-rule principle {[}2,3{]} expressed with
explicit assumptions. It establishes a current-pattern result, not
retention of every earlier pattern, global policy improvement, or
correct interpretation of the target.

\subsection{6.2 An implementable convex feasible
set}\label{an-implementable-convex-feasible-set}

The default graph write changes one provenance tier at a time, with
\(S\) and the other tier fixed. After decay, let \(o_e\) be the other
tier's value, \(p_e^-\) the chosen residual, and
\(z_e^-=S_e+o_e+p_e^-\). First verify the unchanged endpoints
\(|S_e|\le w_{\max}\) and \(|S_e+o_e|\le w_{\max}\) and
\(|o_e|\le p_{\max}\). Define

\[
\begin{aligned}
R_e&=p_{\max}-|o_e|,\\
L_e&=\max(-R_e,-w_{\max}-S_e,-w_{\max}-S_e-o_e),\\
U_e&=\min(R_e, w_{\max}-S_e, w_{\max}-S_e-o_e).
\end{aligned}
\tag{18}
\]

These are the exact interval bounds for the chosen residual under
Equation (4), with the other values fixed. Check
\(L_e\le p_e^-\le U_e\). Allocate nonnegative step allowances \(d_e\)
before the write, with
\(\sum_{e\in\mathcal I_n}d_e\le B_n^{\mathrm{write}}\). Then

\[
\begin{aligned}
l_e&=\max(S_e+o_e+L_e,\ z_e^- -d_e),\\
u_e&=\min(S_e+o_e+U_e,\ z_e^- +d_e),\\
\mathcal C_n&=\prod_{e\in\mathcal I_n}[l_e,u_e].
\end{aligned}
\tag{19}
\]

The post-decay state belongs to this box by construction. Nonselected
coordinates are fixed. Euclidean projection is exactly componentwise
clipping to these intervals. It enforces all four endpoints and the
intentional write bound \(\|z^+-z^-\|_1\le B_n^{\mathrm{write}}\). This
replaces the symmetric residual interval centered at zero used in v0.4;
do not implement that old interval and label it v0.4.1.

Storage-only intervals are obtained by omitting \(z^-\pm d\). They are
used for the reachability diagnostic. Step-limited intervals govern
actual writes. Conflating these sets would mistake a small allowed step
for a lack of representational capacity. Allowance priorities are
bounded allocation heuristics, not usefulness oracles. A zero-width
interval is legal; widespread no-progress projection is a diagnostic,
not evidence of successful learning.

\subsection{6.3 Projected graph write}\label{projected-graph-write}

Use the exact sparse quadratic in Equation (14):

\[
\begin{aligned}
\eta_n&=\frac{\beta_n}{\|A_n\|_F^2+\varepsilon_n},
\qquad 0<\beta_n\le1,\quad\varepsilon_n>0,\\
z^+&=\Pi_{\mathcal C_n}
\big[z^- -\eta_n A_n^\top(A_nz^- -b_n)\big],\\
P_e^{\mathrm{chosen},+}&=z_e^+-S_e-o_e.
\end{aligned}
\tag{20}
\]

Skip the step when \(A_n=0\), the target is unavailable, the record is
inadmissible, or the controller chooses \(\beta_n=0\). For nonzero
\(A_n\), \(L_n=\|A_n\|_2^2\le\|A_n\|_F^2\), so \(\eta_n\le1/L_n\).
Frobenius normalization is conservative and cheap. A certified tighter
bound can improve step size. Disjoint receiver blocks can use separate
step sizes and step allowances when the objective is block-separable. A
decoder that mixes receiver outputs can destroy that separability even
when incoming edge identities do not overlap; then use the joint design
matrix or a justified block method.

\textbf{Proposition 5 --- bounded local descent.} Let
\(\Delta=z^+-z^-\). Under the fixed-feature, fixed-target, fixed-set
assumptions,

\[
\ell_n(z^+)\le\ell_n(z^-)
-\left(\frac1{\eta_n}-\frac{L_n}{2}\right)\|\Delta\|_2^2.
\tag{21}
\]

\emph{Proof.} Projection optimality, using \(z^-\in\mathcal C_n\), gives
\(\langle\nabla\ell_n(z^-),\Delta\rangle\le-\|\Delta\|^2/\eta_n\). The
quadratic smoothness inequality gives
\(\ell_n(z^+)\le\ell_n(z^-)+\langle\nabla\ell_n(z^-),\Delta\rangle+(L_n/2)\|\Delta\|^2\).
Combining them yields Equation (21). Its coefficient is positive for the
reference step size. A zero step may indicate an optimum, an inactive
constraint face, or exhausted capacity; it does not imply successful
task learning.

This theorem applies to the graph's stated local fitting problem.
Unfreezing routing, differentiating through a moving encoder, changing
the target during the step, or using the original weight-dependent
normalizer changes that problem. Such variants require an actual
derivative of their computation and an appropriate step-size or
line-search analysis.

\subsection{6.4 Feasible approximations and finite
precision}\label{feasible-approximations-and-finite-precision}

The reference check suite tests the exact box projection in
floating-point arithmetic against a tolerance. An approximate optimizer
does not inherit the theorem solely because it is called a projection.
If its result satisfies

\[
\langle\nabla\ell_n(z^-),\Delta\rangle
\le-\|\Delta\|^2/\eta_n+\epsilon_{\mathrm{proj}},
\tag{22}
\]

then Equation (21) gains the additive term \(\epsilon_{\mathrm{proj}}\).
A positive predicted descent smaller than numerical or projection error
is not a certified useful improvement. Check the actual captured loss
and feasibility; reject nonfinite results, violated budgets, and
increases beyond declared tolerance. Hot arithmetic starts in FP32 for
correctness experiments; lower precision is a separate ablation.

\subsection{6.5 Local retention
constraints}\label{local-retention-constraints}

A write can fit a new pattern while damaging old patterns. Replay can be
included in a stacked quadratic with fixed nonnegative weights:

\[
\ell_n^{\mathrm{mix}}(z)
=\frac12\sum_{r\in\mathcal B_n}\omega_r
\|A_rz-b_r\|_2^2,
\qquad \omega_r\ge0.
\tag{23}
\]

The same proof applies to the stacked weighted design matrix. It
guarantees descent of the mixture, not of each constituent loss.
Sampling weights define the intended objective. When a different
prioritized sampling distribution is used, either apply an appropriate
importance correction on its supported records or declare that the
optimized distribution has changed.

An optional stronger local preservation set adds constraints such as
\(\|A_r(z-z^-)\|_2\le\epsilon_r\) for protected snapshots. Their
intersection with the box is convex but generally not separable. It
requires an appropriate solver, projection accuracy accounting, and
additional data access. It is not silently implemented by coordinate
clipping. No finite snapshot set establishes preservation on every
possible old input.

\subsection{6.6 Reachability and conditioning before tuning the
learner}\label{reachability-and-conditioning-before-tuning-the-learner}

For diagnostic snapshots, define the storage-constrained error floor

\[
E_{\mathrm{floor}}=\min_{z\in\mathcal C_{\mathrm{storage}}}\|Az-b\|_2^2.
\]

Use the same selected tier, fixed other tier, and slot identities as the
write, but omit the per-step allowance. Record free-coordinate rank,
singular values, an explicitly defined conditioning statistic, active
bound constraints, solver status, and achieved residual. A receiver with
fewer independent writable message directions than target dimensions
cannot fit arbitrary targets. Rank alone is not enough: weight bounds
and target scale also matter.

For an independently checked box solution \(\widehat z\), let
\(y=A\widehat z-b\) and \(c=A^\top y\). For the half-squared objective,
a computable dual lower bound is

\[
D(y)=-\tfrac12\|y\|^2-b^\top y+
\sum_e\min(c_el_e,c_eu_e).
\]

This follows from the Fenchel representation of the squared norm and
minimization of a linear function on each interval. The primal value
\(\tfrac12\|A\widehat z-b\|^2\) is an upper bound on the optimum when
\(\widehat z\) is feasible; \(\max(0,D(y))\) is a lower bound in exact
arithmetic. Report their gap and numerical tolerance. A large achieved
residual without an appropriate lower bound or converged solution does
\textbf{not} establish an irreducible error floor. A bound-constrained
least-squares implementation can use an independently tested numerical
solver {[}13{]}. Fixed coordinates are removed before solving and
reinstated afterward.

Interpretation: a high certified floor suggests deficient features,
connectivity, target coding, or storage capacity; a low floor with weak
actual progress suggests conditioning, step size, or allowances; good
snapshot fit with poor future queries suggests moving-feature or task
misalignment. These are diagnoses to test, not automatic causal
attributions. Full-snapshot oracle fitting is an offline diagnostic and
is never given hidden evaluation labels or used as an undeclared
deployment update.

\subsection{6.7 Measure downstream write
utility}\label{measure-downstream-write-utility}

Clone the complete pre-write state and all random-generator states.
Compare a written and unwritten branch on the same causally valid query
stream:

\[
U_n=\mathcal L_{\mathrm{query}}(\mathcal Z_{\mathrm{no\ write}})
-\mathcal L_{\mathrm{query}}(\mathcal Z_{\mathrm{write}}).
\]

Positive utility indicates better later performance on that stream. Keep
this offline probe from altering live state or revealing probe labels to
future decisions. Report the distribution of \(U_n\), not only its mean.
A locally descending step with negative utility is a concrete warning
that the snapshot target or its features are misaligned with later use.
For short exact-gradient references, also measure alignment between the
proposed update and the negative full query-loss gradient, evaluated
from matched state; alignment is local and is not a long-horizon
guarantee.

\section{7. Interference, capacity, and distinguishable
context}\label{interference-capacity-and-distinguishable-context}

\subsection{7.1 Exact interference in the unconstrained
block}\label{exact-interference-in-the-unconstrained-block}

For an old key \(q\), Equation (16) changes the old prediction by

\[
(W^*-W^-)q
=\frac{\beta e}{\varepsilon+\|k\|_2^2}(k^\top q).
\tag{24}
\]

Orthogonal keys do not interfere in this ideal unconstrained model.
Highly overlapping keys can. The geometry must be learned or engineered;
it is not supplied by the word ``plasticity.'' After projection,
Equation (24) need not hold; use the actual change \(\Delta Wq\) or its
norm bound instead.

\textbf{Restricted existence construction.} Let \(k_1,\ldots,k_m\) be
orthonormal, \(m\le d_k\), and let their targets be within the
unconstrained representable range. Starting at zero, one full delta
write for each key yields \(W=\sum_i v_i k_i^\top\). Every \(Wk_i=v_i\).
A later write for key \(k_j\) changes only its corresponding
association. Exact full transfer \(S\leftarrow S+P\), \(P\leftarrow0\)
preserves all these answers. This establishes coexistence of learning,
retention, selective revision, and consolidation in a restricted memory,
not learned representation discovery from raw observations.

For a counterexample, start with \(W=(1,0)\), use
\(k=(0.9,\sqrt{1-0.9^2})\), and write target \(-1\) with a full delta
update. The new target is fitted, but the old prediction on \((1,0)\)
becomes \(-0.71\). The accompanying tests reproduce both constructions.

\subsection{7.2 Conflicts are not all
forgetting}\label{conflicts-are-not-all-forgetting}

If the same available context and key must simultaneously satisfy
\(Wk=v_1\) and \(Wk=v_2\) with \(v_1\ne v_2\), no deterministic mapping
can satisfy both. Supply a distinguishable task context, entity, time,
or observable history; otherwise mark the earlier relation obsolete or
use a distributional target. Retention metrics must distinguish
still-valid knowledge from intentional correction of outdated knowledge.

The episodic and context systems should store validity intervals, source
identity, and supersession relationships where the task supports them.
In partially observed tasks, the agent needs enough history to infer
context. A benchmark that hides that history cannot attribute every
conflict to poor consolidation.

\subsection{7.3 Finite capacity}\label{finite-capacity}

With \(B\) stored bits there are at most \(2^B\) storage states; more
independent arbitrary information cannot be represented exactly without
collisions. Rewiring reallocates finite capacity; it does not create
unlimited memory. The relevant hypothesis is an improved
adaptation--retention frontier at a declared capacity, precision, and
retrieval budget.

Useful responses to interference include context-sensitive allocation,
rehearsal of representative associations, separation of incompatible
codes, higher-rank or additional blocks, and declared forgetting. All
consume representation capacity, data, or computation. Learned key
diversity should be evaluated jointly with useful generalization: making
every example orthogonal is not a general solution when similar examples
should share structure.

\section{8. Delayed feedback, reinforcement learning, and meta-learning
signals}\label{delayed-feedback-reinforcement-learning-and-meta-learning-signals}

\subsection{8.1 Feedback records instead of unspecified correlation
traces}\label{feedback-records-instead-of-unspecified-correlation-traces}

For supervised or predictive delayed targets, retain the required
fixed-feature record until the target arrives or a declared retention
deadline expires. The writable set can include currently inactive
sources. A fading timestamp alone does not reconstruct missed updates or
identify which earlier decision received the reward.

On a topology or feature-version mismatch, reconstruct the record from
retained admissible observations under the new model, or reject it with
a logged reason. Reusing an old slot's scalar record after its target
changed is semantic corruption. Pending records are pinned or explicitly
spilled; deadlines and overflow cause measured lost-credit events, not
silent erasure.

\subsection{8.2 Reinforcement learning is a separate guarantee
class}\label{reinforcement-learning-is-a-separate-guarantee-class}

A scalar reward or advantage is not inherently inadequate. Under
appropriate policy-gradient assumptions, a sampled update uses a scalar
advantage with parameter-specific sensitivity:

\[
\widehat g_t=\widehat A_t\,
\nabla_\vartheta\log\pi_\vartheta(a_t\mid\mathscr H_t).
\tag{25}
\]

The vector derivative, not the dimensionality of the reward, assigns
different directions to different parameters {[}6{]}. For an adaptive
recurrent agent, exact sensitivities include relevant state and learning
history. Truncated BPTT, local eligibility approximations, and detached
fast-state derivatives must be labeled and compared against small exact
references. E-prop provides an architecture-specific precedent, not a
plug-in proof for this graph {[}5{]}.

Bandit feedback does not reveal the full correct action vector; do not
quietly use the supervised association target in a reward-only
experiment. A value target or auxiliary next-observation target can be
defined, but its regression improvement is not a policy-improvement
theorem. Arbitrary replay is off-policy. Use a justified off-policy
objective, prediction/value fitting, or behavior distillation, and
retain behavior probabilities when an importance-based estimator needs
them.

\subsection{8.3 Inner partial derivatives and outer
meta-gradients}\label{inner-partial-derivatives-and-outer-meta-gradients}

Let \(A_\theta\) denote features produced by shared parameters. The
inner rule computes \(\nabla_z\ell=A_\theta^\top(A_\theta z-b)\) while
holding those features fixed with respect to \(z\). The resulting update
is still a function \(z^+(\theta)\). For a query objective,

\[
\frac{dJ}{d\theta}=\frac{\partial J}{\partial\theta}
+\frac{\partial J}{\partial z^+}\frac{dz^+}{d\theta},
\]

with additional recurrent-state paths included by the full unroll.
Numerical snapshot immutability does not require deleting this second
path. The short-episode reference retains continuous autograd paths
through features, selected gates, message construction, rate
normalization, and the exact functional projection. External target
codes do not require gradients. An intentionally frozen target encoder
stays frozen. The reference does not differentiate through future
evidence that was unavailable at the decision.

Discrete support selection is fixed within a derivative calculation.
Ordinary autodiff gives the exact derivative of the selected smooth
branch away from selection boundaries, not an exact gradient through
changing top-\(k\) membership. State the treatment of routing
explicitly: fixed-support references, smooth training relaxations, score
estimators, and straight-through surrogates are different experiments.
At clip/max/min ties, derivatives can be nonunique. Gradient
finite-difference tests must avoid such kinks, cover interior and stable
active-face cases separately, and report the nondifferentiable cases
rather than claim a false equality.

Use a named \texttt{detached\_feature\_write\_path} ablation to remove
feature-to-write meta-gradients while retaining direct query paths. Do
not accidentally disable all training by running the complete
training-time write under a no-gradient context or converting
graph-connected values to NumPy. Deployment may use detached snapshots
because shared \(\theta\) is then frozen. Truncated episode gradients
and learned local teachers are labeled approximations. Differentiable
plasticity and test-time-training supply relevant precedents for
optimizing how a network learns {[}1,4{]}, not a guarantee that this
training succeeds.

The outer learner may choose representations, continuous route
parameters, allocation, and bounded \(\beta_n\), but cannot override the
local step bound while claiming Proposition 5. A target encoder could
become constant, a router could bypass plastic memory, or a readout
could solve training tasks through unintended context. Task-grounded
losses, write-utility probes, and memory interventions therefore
accompany every claimed learning result. Zero plasticity remains a valid
baseline; activity alone is not evidence of useful learning.

\subsection{8.4 Address stability as an explicit
ablation}\label{address-stability-as-an-explicit-ablation}

Compare the default state-dependent address keys against an
epoch-stable-address variant. In the latter, the address key depends on
stable descriptors and bounded observable cue/context features,
excluding hidden activity produced by the currently writable memory.
Both use the same context information and address dimension. This
isolates feedback in addressing; it does not remove recurrent graph
computation.

A further variant freezes write-message features across equivalent
cue/context queries while the workspace still processes retrieved
information. It is a separate experiment because it changes more than
routing. Record drift between write and later read features, support
overlap, and the relation to write utility. Stable encoders do not by
themselves make state-dependent addresses stable. Retain a variant
because of measured benefit, not because either static or dynamic
routing sounds inherently preferable.

\section{9. Working memory, episodic memory, and
provenance}\label{working-memory-episodic-memory-and-provenance}

\subsection{9.1 Bounded working memory}\label{bounded-working-memory}

Working memory has \(K\) key--value slots, a no-write action, and
explicit retention/eviction metadata. For a selected slot \(a\), use

\[
\begin{aligned}
\widehat v_a&=V_{\max}\tanh(E_v(x_t,h_t)),\\
v_a^+&=(1-\zeta_a)v_a+\zeta_a\widehat v_a,
\qquad0\le\zeta_a\le1,\\
\widehat k_a&=\frac{E_k(x_t,h_t)}{\max(1,\|E_k(x_t,h_t)\|_2)}.
\end{aligned}
\tag{26}
\]

Keys can be replaced or convexly interpolated within the unit ball.
Values remain coordinate-bounded when initialized within that bound.
Read with a norm-bounded query and a bounded similarity plus declared
age bias. A fixed slot count is a storage budget, not a claim that a
slot equals one human memory item.

\subsection{9.2 Episodic record
contract}\label{episodic-record-contract}

A record identifies its source, times, validity/context, observations or
recoverable features, actions, feedback, terminal status, model and
key-encoder versions, and direct admissibility classification. It
contains enough state or history to reconstruct its declared
training/evaluation episode. A short context window is sufficient only
when the task or a fading-memory argument justifies that truncation.

The store has a finite byte budget. Sampling and compaction trade off
coverage, trust, recency, conflict, and task diversity. Synthetic replay
is identified as model-generated and is not promoted to independently
verified evidence by repetition. Keep separate losses for still-valid
retained knowledge and obsolete assertions. Store next observations and
behavior probabilities whenever the chosen objective requires them.

Keys use a stable encoder within each online epoch. Changing its
coordinates requires re-embedding from retained inputs or a validated
compatibility map. Claiming that an index is approximate does not excuse
mismatched feature spaces.

\subsection{9.3 What provenance tiers do and do not
guarantee}\label{what-provenance-tiers-do-and-do-not-guarantee}

\(P^a\) and \(P^u\) separate \textbf{direct write bookkeeping}. The
maintenance controller authorizes which class an event may influence.
Shared nonlinear hidden states can nevertheless transmit indirect
effects of untrusted history into apparently admissible updates. Two
buffers are therefore not a proof of clean attribution, selective
unlearning, or poisoning immunity.

For stronger promotion control, rebuild a slow candidate from a known
checkpoint using independently admissible records and a declared context
policy. Discard or reconstruct histories contaminated beyond that
policy. Provenance judgments remain fallible; independent evidence,
repetition limits, source quotas, and bounded write influence reduce
specified risks but do not prove adversarial robustness.

The protected controller, holdout records, version journal, and action
supervisor are outside the self-modifying set. Protection concerns
access rights and protocol behavior; it does not make their input
signals truthful. External actions taken under a bad fast update cannot
be undone by restoring weights.

\section{10. Consolidation: immediate continuity and durable
recall}\label{consolidation-immediate-continuity-and-durable-recall}

\subsection{10.1 An optional exact
transfer}\label{an-optional-exact-transfer}

For an admissible transfer fraction \(q_e\in[0,1]\), hold topology,
routes, shared parameters, and the other residual tier fixed:

\[
S'_e=S_e+q_eP_e^a,\qquad
P_e^{a\prime}=(1-q_e)P_e^a,\qquad
P_e^{u\prime}=P_e^u.
\tag{27}
\]

\textbf{Proposition 6 --- immediate transfer continuity and endpoint
feasibility.} Equation (27) preserves \(S+P^a+P^u\) exactly in real
arithmetic. The new slow endpoint \(S+qP^a\) lies between the old
endpoints \(S\) and \(S+P^a\). The new slow-plus-ephemeral endpoint
\(S+qP^a+P^u\) lies between \(S+P^u\) and \(S+P^a+P^u\). The other two
endpoints are unchanged. Each remains bounded by Equation (4), and
\((1-q)|P^a|+|P^u|\) cannot increase. Thus exact transfer remains
feasible under the revised endpoint set. This proof does not apply to an
additional quantizer without the separate checks in Section 10.2.

This equality concerns the instantaneous effective function under
unchanged routing and state. It does not preserve the controller's
pressure statistic or all future trajectories. For a single residual
with subsequent decay \(\lambda\) and no new writes, transferring \(qP\)
changes the future effective value relative to no transfer by

\[
qP(1-\lambda^m)
\tag{28}
\]

after \(m\) decay steps. Increased persistence is a purpose of
consolidation, not an algebraic defect. Transfer is optional bookkeeping
and a possible candidate initialization; it is not by itself evidence
that a useful concept has been learned.

\subsection{10.2 Quantization participates in the forward
function}\label{quantization-participates-in-the-forward-function}

Let \(Q_b\) be a declared block quantizer. A compensated transfer is

\[
S'_e=Q_b(S_e+q_eP_e^a),\qquad
P_e^{a\prime}=P_e^a+S_e-S'_e,
\qquad P_e^{u\prime}=P_e^u.
\tag{29}
\]

\textbf{Proposition 7 --- quantized bookkeeping identity.} The sum
remains exactly equal in real arithmetic, including when \(Q_b\)
saturates. Saturation does, however, invalidate a small-error or
bounded-residual promise. Check Equation (4), representable range, and
all storage budgets before accepting this candidate. If infeasible,
change scale or precision, reduce the transfer when that helps, or
reject it. Do not silently clip the compensation.

For a nonsaturating nearest-grid quantizer with spacing \(\Delta_b\),
each rounding error is at most \(\Delta_b/2\). For unbiased stochastic
rounding to adjacent grid points the error magnitude can approach
\(\Delta_b\), with zero conditional mean under the usual in-range
construction. Neither fact proves task accuracy. Quantizer-scale changes
modify other values in the block and must be included in validation.
Low-precision learning research supports testing these mechanisms, not
assuming they work at every precision {[}10{]}.

The compensation must participate in the effective forward weight. An
error accumulator used only in a future optimizer does not preserve the
current forward function. Any cold residual or exception record is
counted as storage. Exact arbitrary real-valued preservation, pure INT4
storage, and removal of every residual cannot all be promised
simultaneously.

\subsection{10.3 Consolidate behavior with fast state
disabled}\label{consolidate-behavior-with-fast-state-disabled}

Let \(M_{S,G,0}\) denote the candidate with both residual tiers zero.
The default sleep fit keeps shared \(\theta\) fixed and optimizes slow
weights on admissible records, followed by limited structural proposals:

\[
\begin{aligned}
J_{\mathrm{sleep}}(S,G)=
&\;\mathbb E_{D_{\mathrm{new}}}\ell(M_{S,G,0},y)\\
&+\lambda_o\mathbb E_{D_{\mathrm{valid\ old}}}\ell(M_{S,G,0},y)\\
&+\lambda_d\mathbb E_{D_{\mathrm{distill}}}
D_{\mathrm{KL}}(p_{\mathrm{ref}}\|p_{S,G,0})
+\lambda_c C(S,G).
\end{aligned}
\tag{30}
\]

All coefficients are nonnegative and declared. The cost term is
optional; it is not a surrogate for measured latency. Distillation uses
still-valid reference behavior and is not treated as ground truth when
that behavior is known to be wrong. The sleep objective is generally
nonconvex through the full graph; Proposition 5 is not a convergence
theorem for Equation (30).

\textbf{Cold-recall protocol.} Disable residual writes and both residual
tiers for the entire evaluation, not just at its first step. Reset or
standardize workspace and working memory. Disable episodic answer
retrieval when measuring pure synaptic retention. Supply only the
context legitimately needed to identify the query. Report a separate
retrieval-enabled condition. Otherwise working memory, residual
reconstruction, or external lookup could conceal failed consolidation.

A function can be consolidated by fitting a slow candidate without
copying its entire transient state. Conversely, copying a scalar can
preserve a transient mistake. The claim tested by H2 is useful cold
recall, not a high fraction of transferred bytes.

\section{11. Maintenance scheduling, validation, and commit
semantics}\label{maintenance-scheduling-validation-and-commit-semantics}

\subsection{11.1 Pressure and bounded maintenance
demand}\label{pressure-and-bounded-maintenance-demand}

Monitor each region's hot population rather than computing a quantile
over predominantly zero cold slots. One normalized load is

\[
s_e=(|P_e^a|+|P_e^u|)/p_{\max},\qquad
S_t=\max_r Q_q\{s_e:e\in\mathcal H_t\cap E_r\}.
\tag{31}
\]

Define an empty region's statistic as zero and monitor the maximum as
well as the quantile. Pending-feedback occupancy, rejected writes,
trusted probe loss, saturation, and time since the last accepted
maintenance also contribute. A weighted pressure score is only an
operational policy; it is not a stability theorem.

The controller has entry/exit hysteresis, minimum spacing, maximum
maintenance work, a maximum permitted write backlog, and a failure
policy. If maintenance repeatedly fails, it can freeze new admissible
writes, retain or spill pinned records within budget, discard declared
ephemeral state, or request external supervision. It must not enter an
unbounded retry loop while losing feedback invisibly. Globally protected
data are not densely scanned at every environmental event; monitor costs
are included in accounting.

\subsection{11.2 Synchronous reference
protocol}\label{synchronous-reference-protocol}

The initial protocol is deliberately synchronous and versioned. Run
maintenance at an episode/reset boundary or another externally safe
pause. In an embodied setting, a fixed supervisor handles actions during
the pause; a model checkpoint cannot rewind the outside world.

\begin{Shaded}
\begin{Highlighting}[]
\NormalTok{MAINTAIN(reference version v)}
\NormalTok{  1. Pause plastic writes; fix the feedback{-}log cutoff.}
\NormalTok{  2. Snapshot S, residual tiers, G, theta/version, scales,}
\NormalTok{     workspace/WM state, pending records, and controller metadata.}
\NormalTok{  3. Build a shadow candidate from admissible records.}
\NormalTok{     Optional transfer must pass its immediate continuity checks.}
\NormalTok{  4. Fit the slow candidate with residuals disabled.}
\NormalTok{  5. Train and evaluate any bounded structural proposal.}
\NormalTok{  6. Check storage, finite arithmetic, graph, and access invariants.}
\NormalTok{  7. Evaluate the fixed candidate on appropriate fresh holdouts,}
\NormalTok{     including cold recall and affected live{-}state conditions.}
\NormalTok{  8. Reconstruct or reset runtime state and reconcile pending}
\NormalTok{     records; prepare a complete new{-}state manifest.}
\NormalTok{  9. If source version is still v and every required check passes,}
\NormalTok{     publish the manifest atomically. Otherwise discard the candidate.}
\NormalTok{ 10. Resume from the published state; do not apply unresolved}
\NormalTok{     old{-}version feedback to new{-}generation slots.}
\end{Highlighting}
\end{Shaded}

Feedback arriving after the snapshot cutoff is buffered with its source
version and accounted for before writes resume. Buffer overflow follows
the declared external-event policy, never silent loss.

The initial algebraic transfer and subsequent functional learning are
different operations. Exact continuity is required for the former. The
latter intentionally changes the model and is judged by its task,
retention, and live-state constraints. Do not add a compensating
residual that silently cancels all useful learned changes merely to
force equality with the old model.

A new model cannot inherit arbitrary old hidden states without
validation. At reset-boundary commits, reset is explicit. Else
reconstruct from a sufficient retained history or use a validated state
adapter. A short recent window is sufficient only with a supporting task
assumption or fading-memory bound. The first protocol freezes shared
\(\theta\) during deployment; broader parameter changes require
correspondingly stronger reconstruction.

An asynchronous extension must log post-snapshot observations and
feedback, replay them under the new model when appropriate, and
revalidate before publication. Raw weight deltas cannot in general be
rebased across changed features or topology. Copy-on-write storage,
journals, crash recovery, and rollback latency are part of that
extension's measured cost.

\subsection{11.3 Three validation levels, with explicit
cost}\label{three-validation-levels-with-explicit-cost}

The reference experiment separates three levels. \textbf{Level 1: every
transaction} runs deterministic invariants plus a preregistered
empirical screen on development probes, labeled empirical and with no
lifetime probability claim. \textbf{Level 2: scheduled independent
evaluations} use held-out episodes and independently trained seeds to
evaluate hypotheses. \textbf{Level 3: optional formal acceptance} uses
the conditions and sample budgets below. A Level 1 pass is never renamed
a Level 3 certificate.

This is a research scheduling choice, not permission to lower safety
requirements in a consequential deployment. If an application requires a
formal acceptance guarantee, a candidate lacking the required evidence
remains inconclusive or rejected. The synchronous simulator reference
has no real-world actuation consequence.

For illustration, \(J=4\), \(\delta=0.01\), and a radius of \(0.01\) in
Equation (32) require \(n\ge133{,}693\) independent episodes. Validation
can dominate learning. Record its examples, wall time, and byte reads. A
tighter valid statistical method is a separately justified change, not
an unrecorded reduction of the radius.

\subsection{11.4 Optional formal
acceptance}\label{optional-formal-acceptance}

Separate replay/training anchors from acceptance data. Let a candidate
be fixed before drawing \(n\) fresh, conditionally independent
evaluation episodes from a declared stratum. Each episode yields a
predeclared bounded loss in \([0,1]\). For paired loss differences
\(d_i=\ell_{\mathrm{candidate},i}-\ell_{\mathrm{reference},i}\in[-1,1]\),
Hoeffding's inequality gives a valid one-sided radius. For \(J\) tested
metrics, a conservative two-sided allocation is

\[
r_\nu=\sqrt{\frac{2\log(2J/\delta_\nu)}{n}}.
\tag{32}
\]

Accept retention only when
\(\bar d_{\mathrm{old}}+r_\nu\le\epsilon_{\mathrm{ret}}\), and new-task
improvement only when \(\bar d_{\mathrm{new}}+r_\nu\le-\eta_{\min}\), in
addition to deterministic constraints. A separately justified binomial
interval can be used for failure rates. Large radii require more data or
an inconclusive decision, not an assumed pass. Report raw errors and
hard failures as well as bounded statistical metrics.

A lifetime error-spending schedule can use

\[
\delta_\nu=\frac{6\delta_{\mathrm{life}}}{\pi^2\nu^2},
\qquad\sum_{\nu=1}^\infty\delta_\nu=\delta_{\mathrm{life}}.
\tag{33}
\]

This supports a union bound \textbf{only if every constituent test
remains valid conditional on prior decisions}. Repeatedly tuning
candidates on the same revealed holdout violates that prerequisite
unless a suitable reusable-holdout or uniform-guarantee mechanism is
supplied {[}9{]}. The reference uses fresh acceptance episodes; if only
a finite repeatedly used data set is available, report empirical
validation without the lifetime statistical claim. Independent episodes,
not correlated time steps from one trajectory, are the sampling units.

These guarantees concern expected loss on specified test distributions.
They do not certify every future behavior, unknown task, poisoning
strategy, or rare catastrophic event.

\section{12. Budgeted structural
adaptation}\label{budgeted-structural-adaptation}

\subsection{12.1 Edge identity and protected
routes}\label{edge-identity-and-protected-routes}

Each slot has a semantic identity
\((\text{source},\text{target},\text{head},\text{generation})\). Reusing
a slot increments its generation and resets every associated residual,
pending credit reference, importance statistic, age, and routing cache
entry. Unresolved feedback is either reconciled before replacement or
retained with its old identity and handled explicitly.

Reserve structural backbone slots and enforce declared regional and
receiver capacities. A backbone preserves possible routes, not
automatically useful signal transmission. Topology checks and temporal
service checks therefore remain separate. No source is forced to prune a
positive fraction of useful edges.

\subsection{12.2 Candidate selection and actual
value}\label{candidate-selection-and-actual-value}

Use bounded candidate sets from local regions, existing indexed
neighbors, two-hop relations, and random exploration. Choose both target
and message head. Filter forbidden duplicates, self-relations where
disallowed, and capacity violations before sampling. Empty admissible
sets mean ``no change,'' not an undefined softmax. Candidate
construction and any gradient proposal cost are charged.

First-order saliency can screen proposals but cannot establish
indispensability. At \(\ell(w)=\tfrac12(w-1)^2\), the correct \(w=1\)
has \(|w\ell'(w)|=0\) although pruning it incurs loss \(1/2\). Estimate
the loss of actual ablation and validate the full network, including
route normalization. Magnitude/gradient methods such as RigL provide
important comparison points {[}7{]}; CSGN's proposed search is not
assumed superior.

\subsection{12.3 Admission and training of new
edges}\label{admission-and-training-of-new-edges}

A new live slot begins unavailable and with zero semantic state. Its
zero-availability behavior follows Equation (7). However, zero
availability also gives zero functional gradient through that edge. To
learn a useful candidate, enable it with positive availability on the
\textbf{shadow model}, assign probationary routing capacity, and train
it on admissible data before considering publication.

Availability may ramp on a declared schedule, with age advancing
independently of successful usage. Reserve probationary support or
otherwise account for incumbent displacement. A positive-availability
support change is not claimed function-preserving; evaluate the whole
candidate. The reference transaction budget caps replaced slots,
training steps, data bytes, and affected regions. Rewiring is retained
only if it improves H3 beyond matched fixed-topology and standard
sparse-training baselines.

\section{13. Complete-lifecycle training and
evaluation}\label{complete-lifecycle-training-and-evaluation}

\subsection{13.1 Training phases}\label{training-phases}

The first phase learns bounded interface features, routing, graph
computation, and readout on task-grounded examples. The second enables
the projected local writes while shared parameters are trained through
manageable episode lengths. The third trains and evaluates sleep
followed by residual clearing and cold recall. Structural adaptation is
introduced after that lifecycle is functional, then included in
integrated tests and ablations.

For a sampled task sequence \(\tau_{1:T}\), the intended outer criterion
is

\[
\min_{\theta,\phi}\;
\mathbb E_{\tau_{1:T}}
\sum_{t=1}^T\left[
\ell_{\mathrm{query},t}
+\lambda_{\mathrm{cold}}\ell_{\mathrm{cold},t}
+\lambda_{\mathrm{old}}\ell_{\mathrm{valid\ old},t}
+\lambda_{\mathrm{cost}}C_t
\right].
\tag{34}
\]

\(\phi\) denotes bounded write/allocation and maintenance-policy
parameters. The state evolves through the declared compute, write, and
maintenance operations. Not every transaction decision is
differentiable. Fixed controller choices, outer-loop search, surrogate
gradients, and truncated differentiation are labeled by experiment; they
are not presented as exact differentiation through an unlimited
lifetime. Only locally guaranteed updates retain their analytical step
bounds.

Training sequences explicitly include learning, holding, distractors,
delayed feedback, maintenance, residual clearing, recall, context
changes, and relearning. Hold out cue identities, context combinations,
and whole task families as appropriate. Evaluate horizons and delay
distributions longer than those used in meta-training.

\subsection{13.2 Small reference profile}\label{small-reference-profile}

The following fixes the size and scheduling profile of a first
experiment, not a validated hyperparameter optimum. Its task-specific
encoder and value-code interface must also be selected and recorded:

\begin{longtable}[]{@{}
  >{\raggedright\arraybackslash}p{(\columnwidth - 2\tabcolsep) * \real{0.5000}}
  >{\raggedright\arraybackslash}p{(\columnwidth - 2\tabcolsep) * \real{0.5000}}@{}}
\toprule\noalign{}
\begin{minipage}[b]{\linewidth}\raggedright
Quantity
\end{minipage} & \begin{minipage}[b]{\linewidth}\raggedright
Reference choice
\end{minipage} \\
\midrule\noalign{}
\endhead
\bottomrule\noalign{}
\endlastfoot
Columns and regions & \(N=512\); 16 regions of 32 columns \\
Reserved slots & \(k_{\mathrm{out}}=16\); \(E=8,192\) \\
Slot classes & 4 protected slots and 12 plastic slots per source \\
State and messages & \(d=64\), \(d_m=16\), \(H=4\) \\
Active budget & At most 4 selected regions, 64 senders, 4 real executed
routes per sender \\
Inner rounds & \(K_{\mathrm{inner}}=4\) with frozen route support \\
Workspace regime & Sufficient norm condition in Equation (12) \\
Online parameters & Shared \(\theta\) fixed; projected writes to
selected residual tier \\
Initial precision & FP32 slow and hot state; float64 isolated algebra
checks \\
Association targets & Bounded externally supplied value codes, revealed
after prediction \\
Retention test & Admissible residual decay disabled; ephemeral decay
separately declared \\
Maintenance & Synchronous episode-boundary candidate fitting and
commit \\
Rewiring & Fixed in the first baseline; bounded shadow proposals in
structural experiments \\
\end{longtable}

Choose weight/rate/data budgets before the run and publish them in its
manifest. Message targets must lie in a sensible scale relative to the
receiver bound; target reachability is measured, not assumed. Hot
capacity, pending-feedback lifetime, replay bytes, and acceptance sample
budgets are explicitly swept rather than left unbounded. Pilot defaults
in the handoff configs are engineering starting points, not empirical
recommendations. Keep the full compute--write--sleep--clear--recall
lifecycle. Defer rewiring, online shared-encoder changes, INT4,
offloading, asynchronous maintenance, and learned internal teachers
until their own prerequisites pass. Reference region scoring is exact
and its cost is reported. Subsequent scales can use \(N=4,096\) and then
larger systems only after evidence warrants them.

\subsection{13.3 Required comparisons}\label{required-comparisons}

Use a static recurrent graph with the same encoders and training budget;
a replay-only slow learner; the prior bounded correlation rule as an
ablation; normalized delta/block-memory baselines; the v0.4.1 graph
writes; and v0.4.1 with/without consolidation and rewiring. Established
fast-memory results motivate these comparisons {[}1--4{]}, not a claim
that a particular baseline must lose.

Match available observations, label timing, task identifiers,
representation training, total parameter bytes, fast-state bytes, replay
exposure, inference rounds, and wall-clock or operation budgets. Give a
replay-only baseline the replay compute that sleep consumes. Report both
matched-memory and matched-latency comparisons when simultaneous
matching is not practical. Test-time interventions that disable writes
complement, but do not replace, separately trained baselines.

\subsection{13.4 Causal memory interventions and
transfer}\label{causal-memory-interventions-and-transfer}

\textbf{Slow-memory transplantation.} Use two agents with identical
shared \(\theta\), descriptors, readout, topology, initial state, and
quantization convention, but different randomly assigned valid
association tables. After separate learning/consolidation, transplant
only slow synaptic values from donor to recipient. Reset both residuals
and all transient state, disable episodic answer retrieval and plastic
writes, and evaluate the donor's legitimate cue/context queries.
Recalled associations should track the transplanted memory if the
demonstrated knowledge resides there. Do not transfer a decoder or
task-ID lookup with the weights and call that synaptic causality. Use
fixed topology initially so slot identities are compatible. Independent
training of \(\theta\) would invalidate this simple intervention.

\textbf{Selective erasure.} Select candidate memory groups using
training/development data only; compare their erasure with equal-size
and magnitude/usage-matched random erasure. Test affected and unaffected
query classes with uncertainty. Generic model damage is not selective
causal localization. In this experiment erasure is an offline
intervention, not a silently deployed structural edit.

\textbf{Negative controls.} Fresh independent labels on every query,
with no predictive information in any permitted history, should not
yield held-out prediction above the declared chance expectation. Stable
randomly assigned cue labels are \emph{not} a nonlearning control: they
can legitimately be memorized after feedback. Also randomize support
feedback independently of scored query truth to detect leakage, and
assert that future targets, task-generator seeds, and hidden map tables
never enter observations.

\textbf{Rule transfer.} Include episodes whose feedback identifies a
structured context-dependent transformation, then query unseen inputs
and revisit the rule after distractors and sleep. The support set must
identify the chosen family; do not demand arbitrary linear-map recovery
from insufficient independent examples. For example, a
signed-permutation linear map can be tested after a full-rank support
design. Hold out input vectors, mappings, and, separately, map families.
Shared-representation meta-learning theory motivates specifying such
structure {[}12{]}; it does not certify this nonlinear graph.

The simple delta-memory comparator is a primary competitor on these
tasks. Match information and training budgets; compare matched-memory
and matched-latency frontiers when exact simultaneous matching is
impractical. Report uncertainty across independently trained seeds, not
only resampled queries from one checkpoint {[}11{]}. Any transplant or
transfer claim requires the integrated trained graph, not the ideal
matrix checks in Section 15.

\subsection{13.5 Metrics and decision
gates}\label{metrics-and-decision-gates}

Measure current-task error/regret, few-shot adaptation, still-valid
old-task loss, cold recall, context-specific revision, retrieval-enabled
recall, and generalization to held-out tasks. Add downstream write
utility, snapshot floor bounds and solver gaps, active-feature
rank/conditioning, address drift, transplantation tracking,
selective-erasure specificity, and held-out rule transfer. Also report
projection saturation, inactive/no-progress writes, interference versus
key similarity, drift across maintenance, rollback/rejection rates,
missed-credit events, routing coverage/latency, and actual bytes and
time.

\textbf{Gate A: local mathematics.} Invariants, current-snapshot
descent, zero-availability behavior, transfer identities, and numerical
tolerances pass their dedicated checks.

\textbf{Gate B: learned adaptation.} Learned features and routing
produce reproducible gains beyond static and replay-only models; success
with externally orthogonal keys is not sufficient.

\textbf{Gate C: durable consolidation.} Pure slow recall survives
residual disabling and transient-state reset on both new and still-valid
prior tasks.

\textbf{Gate D: delayed and changing context.} The agent handles delayed
targets, distractors, revised facts, and longer horizons without hidden
target leakage or impossible identical-context requirements.

\textbf{Gate E: structural value and cost.} Rewiring improves an
adaptation--retention or resource frontier after its overhead is
charged.

\textbf{Gate F: robustness and scale.} Adversarial, noisy,
distribution-shift, precision, and resource-pressure tests do not reveal
unacceptable failures, and measured sparse execution is competitive.

A failed gate prompts diagnosis or simplification. A passed gate
supports only its tested conditions. Preregister the primary metrics and
selection rule; report multiple independent seeds and uncertainty, not
only the best run.

\section{14. Storage and compute
accounting}\label{storage-and-compute-accounting}

\subsection{14.1 Count the entire adaptive
state}\label{count-the-entire-adaptive-state}

Let all \(b_\cdot\) variables denote \textbf{bytes}, not bits. A minimum
accounting decomposition is

\[
\begin{aligned}
M_{\mathrm{total}}={}&E(b_S+b_{\mathrm{target}}+b_{\mathrm{type}}+b_{\mathrm{meta}})\\
&+|\mathcal H|(b_{P^a}+b_{P^u}+b_{\mathrm{hot\ index}}+b_{\mathrm{hot\ meta}})\\
&+Ndb_h+M_{\mathrm{pending}}+M_{\mathrm{WM}}+M_{\mathrm{episodic}}\\
&+M_\theta+M_{\mathrm{routing}}+M_{\mathrm{training}}
+M_{\mathrm{transaction}}+M_{\mathrm{communication}}.
\end{aligned}
\tag{35}
\]

Training memory includes gradients, optimizer state, and unrolled
activations. Transaction memory includes snapshots, candidate
differences, validation, scales, journals, and state migration. Pending
records include retained feature matrices or the data and versions
needed to reconstruct them. Do not hide a full-precision residual or
optimizer value for every cold edge inside an unspecified constant.

Cold INT4 coefficients alone occupy \(E/2\) bytes when perfectly packed.
At \(E=100\) billion this is 50 GB in decimal units; it says nothing
about total residency. Indices, node state, and active histories can
dominate. Node count is \(N=E/k_{\mathrm{out}}\); target-identifier
width must support that actual \(N\). The 86--100B figure remains a
possible distributed storage research target, not a first model
configuration or a biological equivalence.

\subsection{14.2 Pending credit creates a workload-dependent lower
bound}\label{pending-credit-creates-a-workload-dependent-lower-bound}

If \(m\) newly distinct edges per step must remain individually eligible
for \(T\) steps, the union can reach \(\min(E,mT)\) edge identities.
Multiple event-specific snapshots can require more than one record per
edge. A tiny fixed hot fraction cannot simultaneously promise arbitrary
delayed feedback, no information loss, and no external spill. Cap delays
or arrivals, spill within budget, compress with justified algebra, or
declare and measure expired-credit losses.

\subsection{14.3 Scored, executed, and written work
differ}\label{scored-executed-and-written-work-differ}

Let \(\mathcal Q_t\) contain scored edge candidates, \(\mathcal X_t\)
executed edges, \(\mathcal U_t\) updated nodes, and \(\mathcal I_n\) the
write set. A representative cost ledger is

\[
\begin{aligned}
C_{\mathrm{decision}}={}&C_{\mathrm{route\ lookup}}+C_{\mathrm{candidate\ scores}}(|\mathcal Q_t|)\\
&+K_{\mathrm{inner}}\big[
C_{\mathrm{messages}}+O(|\mathcal X_t|d_m)+C_{\mathrm{workspace}}\big]\\
&+C_{\mathrm{readout}}+C_{\mathrm{snapshot}}+C_{\mathrm{queues}}+C_{\mathrm{data\ movement}}.
\end{aligned}
\tag{36}
\]

For a captured sparse quadratic, the gradient costs sparse products by
\(A_n\) and \(A_n^\top\); box projection is \(O(|\mathcal I_n|)\).
Replay and multiple write iterations multiply these costs. Scoring
top-\(k_g\) routes can require inspecting all stored outgoing candidates
even though only \(k_g\) execute. Decoder work, source-head computation,
and index refresh must not be omitted.

Sleep is reported separately in replay examples, optimization steps,
structural candidates, validation episodes, bytes transferred, peak
storage, and elapsed time. With a fixed active fraction, work grows with
graph size. Constant active work requires a decreasing fraction and a
routing/retrieval mechanism that still finds useful information. Count
cache misses, scatter contention, effective bandwidth, and cold-fetch
latency; sparse arithmetic counts alone do not establish a systems
advantage.

\section{15. Verification results and
limits}\label{verification-results-and-limits}

\subsection{15.1 Executed checks and evidence
classes}\label{executed-checks-and-evidence-classes}

The companion \nolinkurl{verification/check_math_v0_4_1.py} was newly
executed for this release on CPU with deterministic seed 20260916. Its
24 groups contain 19,438 local cases using NumPy, SciPy, and PyTorch.
All passed their programmed tolerances. These are numerical reference
and gradient checks, not an integrated CSGN implementation, trained
agent, or experimental result for H1--H3.

\begin{longtable}[]{@{}
  >{\raggedright\arraybackslash}p{(\columnwidth - 4\tabcolsep) * \real{0.3000}}
  >{\raggedleft\arraybackslash}p{(\columnwidth - 4\tabcolsep) * \real{0.4000}}
  >{\raggedright\arraybackslash}p{(\columnwidth - 4\tabcolsep) * \real{0.3000}}@{}}
\toprule\noalign{}
\begin{minipage}[b]{\linewidth}\raggedright
Check family
\end{minipage} & \begin{minipage}[b]{\linewidth}\raggedleft
Executed cases
\end{minipage} & \begin{minipage}[b]{\linewidth}\raggedright
Scope
\end{minipage} \\
\midrule\noalign{}
\endhead
\bottomrule\noalign{}
\endlastfoot
Revised endpoint feasibility, decay and interval equivalence & 6,001 &
Includes saturated slow-weight inward correction \\
Projected writes under revised endpoints & 3,000 & Descent inequality,
endpoint bounds and write budget \\
Normalized delta identity & 5,000 & Fixed keys and targets \\
Exact and quantized transfer checks & 3,001 & Feasibility for exact
transfer; rejection example for quantization \\
Receiver and frozen-workspace examples & 2,000 & Local boundedness and
sufficient contraction examples \\
Reachability references & 103 & Known floors, fixed coordinates, and
primal--dual gaps \\
Outer-gradient and cross-framework checks & 302 & Finite differences,
detached-path ablation, active-face test, NumPy/PyTorch parity \\
Other identities and counterexamples & 31 & Ideal matrix interventions,
interference, routing, invalid inputs, statistical-cost arithmetic \\
\end{longtable}

Raw results, tolerances, interpreter/library versions, timestamps, and
SHA-256 hashes of the actual test sources are retained in
\nolinkurl{verification/results_v0_4_1.json}. The verification script
exits unsuccessfully on a failed assertion. It must not replace failed
tests with cached output. Historical v0.4 results are included only in
the labeled archive.

The ideal matrix transplantation check does not establish memory
transplantation in a trained graph. The gradient check differentiates a
small continuous fixed-support write, not discrete router selection or
an entire learning lifetime. The diagnostic solver's bounds concern a
fixed box-constrained snapshot, not downstream policy capacity. Proofs
and their assumptions remain the basis of analytical claims; random
trials detect inconsistencies in numerical implementations.

\subsection{15.2 Guarantee ledger}\label{guarantee-ledger}

\begin{longtable}[]{@{}
  >{\raggedright\arraybackslash}p{(\columnwidth - 4\tabcolsep) * \real{0.3333}}
  >{\raggedright\arraybackslash}p{(\columnwidth - 4\tabcolsep) * \real{0.3333}}
  >{\raggedright\arraybackslash}p{(\columnwidth - 4\tabcolsep) * \real{0.3333}}@{}}
\toprule\noalign{}
\begin{minipage}[b]{\linewidth}\raggedright
Statement
\end{minipage} & \begin{minipage}[b]{\linewidth}\raggedright
Conditions
\end{minipage} & \begin{minipage}[b]{\linewidth}\raggedright
Not established
\end{minipage} \\
\midrule\noalign{}
\endhead
\bottomrule\noalign{}
\endlastfoot
Receiver messages remain bounded & Fixed nonnegative route factors,
bounded messages/weights, Equation (9) & Informative messages or good
routing \\
Workspace box is invariant & Bounded initial states and convex update &
Useful computation \\
Workspace contracts during a cycle & Frozen conditions and Equation (12)
& Global adaptive or environmental stability \\
Delta write reduces current association error & Fixed key/target, valid
rate, no binding projection & Retention of other keys \\
Projected graph write reduces captured loss & Fixed quadratic, exact
feasible projection, valid step & Improvement of a moving-feature or
final policy loss \\
Transfer preserves the immediate effective sum & Stated bookkeeping and
unchanged forward dependencies & Future trajectory equality or durable
correctness \\
Quantized transfer preserves the sum with compensation & Compensation
remains in forward state; accepted feasibility & Small residual under
saturation or pure cold exactness \\
\end{longtable}

\subsection{15.3 Principal unresolved
risks}\label{principal-unresolved-risks}

The largest scientific uncertainty is whether learned local targets and
feature geometry make snapshot descent useful for downstream tasks.
Other risks include representational collapse, context aliasing,
plasticity suppression, excessive interference, stale features, lost
delayed credit, and constraints that eliminate capacity before useful
learning occurs.

Sleep may fail to fit slow memory, preserve outdated answers, or reject
nearly every useful candidate. Routing can leave relevant regions
dormant or make computation too shallow. Rewiring can consume validation
effort without benefit. Protected controllers can still ingest
misleading evidence. Quantization, replay, and sparse data movement can
erase any practical advantage. These are empirical questions and
specified threat-model questions, not issues solved by adding more
bounded scalars.

\section{16. Conclusion}\label{conclusion}

Version 0.4.1 retains CSGN's sparse graph, fast and slow synaptic state,
working and episodic memory, sleep, and structural plasticity, but gives
them a clearer division of labor. The graph computes with frozen routing
and memory during a bounded cycle. An evidence-backed target defines a
local error. A projected update corrects that error within exact
endpoint-feasible reserves and write budgets, while continuous outer
meta-gradients teach which writes are useful. Sleep fits and evaluates
retained behavior with fast state disabled. Structural edits are
admitted only after their whole-network effects and costs are tested.

The result is a research specification with meaningful local guarantees
and explicit global uncertainties. The v0.4.1 handoff makes those
uncertainties testable: determine target reachability before tuning,
compare write and no-write continuations, transplant slow memory under
controlled state, test rule transfer, and confront the simple-memory
baseline. It is not a theorem that CSGN will become a capable lifelong
agent. The decisive next evidence is a trained small graph whose learned
representations support useful adaptation, whose sleep produces genuine
cold recall, and whose performance survives matched baselines,
interference, delay, and resource pressure.

\section{Appendix A. Focused migration from
v0.4}\label{appendix-a.-focused-migration-from-v0.4}

\begin{longtable}[]{@{}
  >{\raggedright\arraybackslash}p{(\columnwidth - 2\tabcolsep) * \real{0.5000}}
  >{\raggedright\arraybackslash}p{(\columnwidth - 2\tabcolsep) * \real{0.5000}}@{}}
\toprule\noalign{}
\begin{minipage}[b]{\linewidth}\raggedright
Area
\end{minipage} & \begin{minipage}[b]{\linewidth}\raggedright
Required v0.4.1 change
\end{minipage} \\
\midrule\noalign{}
\endhead
\bottomrule\noalign{}
\endlastfoot
Feasibility, Equations (4), (18), (19) & Four-endpoint bounds, exact
asymmetric residual intervals, unchanged separable projection \\
Transfer, Proposition 6 & Replace absolute-sum proof with
endpoint-convexity proof \\
Snapshot and meta-training & Fix the inner variable while retaining
outer continuous gradient paths; explicit detach ablation \\
Local diagnostics & Storage-only reachable floor, solver gap, rank,
conditioning, and future write utility \\
Addressing & Compare state-dependent and stable-address variants without
changing context information \\
Memory evidence & Transplantation, matched erasure, negative controls
and unseen-input rule transfer \\
Baselines & Simple delta memory is a primary competitor; independent
training seeds and fair resource accounting \\
Validation & Empirical per-transaction screen, independent research
evaluations, optional formal certificates \\
Scope & Keep complete small lifecycle; defer complexity without
declaring it unnecessary forever \\
\end{longtable}

Equation numbers (1)--(36) are retained for implementation traceability;
new diagnostics use unnumbered definitions. v0.4 code or constraints are
not alternate defaults. Any intentional deviation is recorded against
the contract map in the handoff and evaluated as a named variant.

\section{Appendix B. Historical migration from
v0.3}\label{appendix-b.-historical-migration-from-v0.3}

\begin{longtable}[]{@{}
  >{\raggedright\arraybackslash}p{(\columnwidth - 2\tabcolsep) * \real{0.5000}}
  >{\raggedright\arraybackslash}p{(\columnwidth - 2\tabcolsep) * \real{0.5000}}@{}}
\toprule\noalign{}
\begin{minipage}[b]{\linewidth}\raggedright
v0.3 area
\end{minipage} & \begin{minipage}[b]{\linewidth}\raggedright
v0.4 disposition
\end{minipage} \\
\midrule\noalign{}
\endhead
\bottomrule\noalign{}
\endlastfoot
Equations 1--5: graph and activation & Retain fixed slot accounting;
specify exact small-scale routing and count search overhead \\
Equation 6: additive context score & Replace with query--key interaction
because the shared additive term cancels \\
Equations 7--17: messages and gates & Preserve bounded typed messages;
include availability in gate selection; use a weight-independent
receiver envelope \\
Equations 18--22: recurrent cell & Separate frozen-cycle computation;
provide a sufficient workspace contraction condition \\
Equations 23--33: correlation/eligibility/fast updates & Replace the
default supervised write with a fixed-objective projected update;
separate retention and writing; specify delayed records \\
Equations 34--36: importance/metaplasticity & Treat scores as allocation
heuristics; bound write rates; require ablation before structural
deletion \\
Equations 37--43: memory and pressure & Bound memory endpoints; use
versioned records, hot-set statistics, explicit overflow and maintenance
policies \\
Equations 44--49: consolidation and acceptance & Distinguish
instantaneous continuity from persistence; include forward quantization
compensation, cold recall, state publication, and valid testing \\
Equations 50--54: structural change & Preserve bounded candidate search;
train unavailable new edges on a shadow model before positive live
admission \\
Equation 55: meta-objective & Evaluate complete
write--hold--sleep--clear--recall lifecycles, not only wake query
loss \\
Equations 56--59: storage and cost & Count both residual tiers,
snapshots, pending histories, node state, scored candidates, and
transaction overhead \\
\end{longtable}

Several prior audit statements are deliberately narrowed. A scalar
reward is not intrinsically inadequate; local correlations are not
automatically gradient eligibilities. Two provenance buffers are not
complete isolation. Summable test error allowances do not repair an
adaptively reused holdout. Strict contraction of every memory coordinate
would defeat persistent storage. Instantaneous transfer is not
future-trajectory equivalence. Each correction is incorporated into the
main specification rather than left as an external caveat.

\section{Appendix C. Reproducibility and
provenance}\label{appendix-c.-reproducibility-and-provenance}

The immediate baseline is the conversation attachment
\texttt{csgn\_paper\_v0\_4.md}; its SHA-256 is recorded in
\texttt{provenance/sources.json}. Historical design lineage includes
\nolinkurl{Patwrick/Neuroplastic/Docs/csgn_paper_v0_3.md}, Git blob
\nolinkurl{dddf9f7a196bc77270dd10415ba7bb07fb4a4805}, and the
objective-based reassessment of 16 September 2026. The package includes
the original v0.4 research archive and reassessment archive unchanged
and clearly non-authoritative.

The repository metadata and README were read at commit
\nolinkurl{281746c9ae962634df04ec34b3ba212af11b3cfd} while preparing
this handoff. That observed commit is a discovery reference, not a
requirement to reset a newer checkout. No repository files were modified
or published during package preparation. Codex must inspect the live
checkout, existing instructions, uncommitted work, dependencies, and
available hardware before implementation.

The package contains this Markdown manuscript, a typeset PDF and LaTeX
source, mathematical reference functions, newly executed checks and raw
results, implementation contracts, task protocols, experiment
configurations, report schemas, a staged Codex prompt, and a revision
feedback loop. The reference functions are not a complete CSGN. There is
no trained checkpoint, actual learning benchmark, or claimed empirical
improvement in this release.

From the extracted package root, use the pinned reference dependency
list to reproduce the CPU checks in an isolated environment:

\begin{Shaded}
\begin{Highlighting}[]
\ExtensionTok{python} \AttributeTok{{-}m}\NormalTok{ pip install }\AttributeTok{{-}r}\NormalTok{ requirements{-}reference.txt}
\ExtensionTok{python}\NormalTok{ verification/check\_math\_v0\_4\_1.py }\AttributeTok{{-}{-}output}\NormalTok{ verification/rerun.json}
\ExtensionTok{python}\NormalTok{ scripts/validate\_package.py}
\end{Highlighting}
\end{Shaded}

PyTorch builds and CUDA installation for the later implementation are
hardware-dependent. Codex must record its actual environment and use
appropriate official installation guidance; old repository GPU scripts
are historical, not a tested requirement of this release. GPU-less
environments can run all supplied reference checks. Do not silently skip
a gradient check and report full verification.

A fresh rerun produces a new result file. Keep the shipped result
immutable as release evidence. BLAS and package differences can change
last-bit values; declared tolerances and source hashes make those
differences inspectable. Package hashes establish byte identity, not
correctness of the scientific claims. The archive should contain no
credentials, external font files, private data, or undisclosed model
outputs.

\Needspace{10\baselineskip}

\section*{References}\label{references}
\addcontentsline{toc}{section}{References}

\small

\textbf{{[}1{]}} Miconi, T.; Stanley, K.; Clune, J. (2018).
\emph{Differentiable plasticity: training plastic neural networks with
backpropagation.} ICML, PMLR 80, 3559--3568.
\url{https://proceedings.mlr.press/v80/miconi18a.html}

\textbf{{[}2{]}} Schlag, I.; Irie, K.; Schmidhuber, J. (2021).
\emph{Linear Transformers Are Secretly Fast Weight Programmers.} ICML,
PMLR 139, 9355--9366.
\url{https://proceedings.mlr.press/v139/schlag21a.html}

\textbf{{[}3{]}} Yang, S.; Kautz, J.; Hatamizadeh, A. (2025).
\emph{Gated Delta Networks: Improving Mamba2 with Delta Rule.} ICLR;
arXiv 2412.06464v3. \url{https://arxiv.org/abs/2412.06464}

\textbf{{[}4{]}} Sun, Y. et al.~(2025). \emph{Learning to (Learn at Test
Time): RNNs with Expressive Hidden States.} ICML, PMLR 267,
57503--57522. \url{https://proceedings.mlr.press/v267/sun25h.html}

\textbf{{[}5{]}} Bellec, G.; Scherr, F.; Subramoney, A.; Hajek, E.;
Salaj, D.; Legenstein, R.; Maass, W. (2020). \emph{A solution to the
learning dilemma for recurrent networks of spiking neurons.} Nature
Communications 11, 3625.
\url{https://doi.org/10.1038/s41467-020-17236-y}

\textbf{{[}6{]}} Sutton, R. S.; McAllester, D. A.; Singh, S. P.;
Mansour, Y. (1999). \emph{Policy Gradient Methods for Reinforcement
Learning with Function Approximation.} Advances in Neural Information
Processing Systems 12.
\url{https://papers.nips.cc/paper/1999/hash/464d828b85b0bed98e80ade0a5c43b0f-Abstract.html}

\textbf{{[}7{]}} Evci, U.; Gale, T.; Menick, J.; Castro, P. S.; Elsen,
E. (2020). \emph{Rigging the Lottery: Making All Tickets Winners.} ICML,
PMLR 119, 2943--2952.
\url{https://proceedings.mlr.press/v119/evci20a.html}

\textbf{{[}8{]}} Zenke, F.; Poole, B.; Ganguli, S. (2017).
\emph{Continual Learning Through Synaptic Intelligence.} ICML, PMLR 70,
3987--3995. \url{https://proceedings.mlr.press/v70/zenke17a.html}

\textbf{{[}9{]}} Dwork, C.; Feldman, V.; Hardt, M.; Pitassi, T.;
Reingold, O.; Roth, A. (2015). \emph{The reusable holdout: Preserving
validity in adaptive data analysis.} Science 349, 636--638. Author
research record:
\url{https://research.ibm.com/publications/the-reusable-holdout-preserving-validity-in-adaptive-data-analysis}

\textbf{{[}10{]}} Gupta, S.; Agrawal, A.; Gopalakrishnan, K.; Narayanan,
P. (2015). \emph{Deep Learning with Limited Numerical Precision.} ICML,
PMLR 37, 1737--1746.
\url{https://proceedings.mlr.press/v37/gupta15.html}

\textbf{{[}11{]}} Agarwal, R.; Schwarzer, M.; Castro, P. S.; Courville,
A.; Bellemare, M. G. (2021). \emph{Deep Reinforcement Learning at the
Edge of the Statistical Precipice.} NeurIPS 34.
\url{https://proceedings.neurips.cc/paper/2021/hash/f514cec81cb148559cf475e7426eed5e-Abstract.html}

\textbf{{[}12{]}} Tripuraneni, N.; Jin, C.; Jordan, M. (2021).
\emph{Provable Meta-Learning of Linear Representations.} ICML, PMLR 139,
10434--10443.
\url{https://proceedings.mlr.press/v139/tripuraneni21a.html}

\textbf{{[}13{]}} SciPy developers. \emph{scipy.optimize.lsq\_linear ---
bounded linear least squares.} Official documentation, accessed 16
September 2026.
\url{https://docs.scipy.org/doc/scipy/reference/generated/scipy.optimize.lsq_linear.html}

\end{document}
